% Learning a Decision Policy from Synthetic Data:
% Pre-Registered QLoRA Fine-Tuning of Open-Weight Models for a Student AI Operating System
%
% Compile with pdfLaTeX + BibTeX (Overleaf: default compiler, main document main.tex).
% Every number in this paper is copied from a record in the KLEOS Models repository;
% the "% source:" comments in the section files name the record and line.

\documentclass[10pt,letterpaper]{article}
\PassOptionsToPackage{hyphens}{url}

\usepackage[letterpaper,margin=1in]{geometry}
\usepackage[T1]{fontenc}
\usepackage[utf8]{inputenc}
\usepackage{newtxtext,newtxmath}
\usepackage{microtype}
\usepackage{amsmath,amssymb}
\usepackage{booktabs,array}
\usepackage[font=small,labelfont=bf]{caption}
\usepackage{enumitem}
\usepackage{xcolor}
\usepackage{tikz}
\usetikzlibrary{positioning,arrows.meta}
\usepackage{pgfplots}
\pgfplotsset{compat=1.18}
\usepackage{natbib}
\usepackage{float}
\usepackage{etoolbox}
\apptocmd{\thebibliography}{\sloppy}{}{}
\usepackage[colorlinks=true,linkcolor=blue!50!black,citecolor=blue!50!black,urlcolor=blue!50!black]{hyperref}

\setlength{\bibsep}{3pt}
\setlist{nosep,leftmargin=*}
\setlength{\parskip}{0pt}
\setlength{\tabcolsep}{3pt}
\newcolumntype{L}[1]{>{\raggedright\arraybackslash}p{#1}}
\newcolumntype{R}[1]{>{\raggedleft\arraybackslash}p{#1}}

% Project names and one notation macro.
\newcommand{\kleos}{KLEOS}
\newcommand{\hermes}{KLEOS Hermes}
\newcommand{\logos}{KLEOS Logos}
\newcommand{\ci}[2]{[$#1$, $#2$]}
\newcommand{\code}[1]{\texttt{#1}}

\title{\Large\bfseries Learning a Decision Policy from Synthetic Data:\\
Pre-Registered QLoRA Fine-Tuning of Open-Weight Models\\
for a Student AI Operating System}
\author{Tejas Naik\\
The Ohio State University\\
\texttt{kleos2027@gmail.com}}
\date{October 2026 (version 1)}

\begin{document}

\maketitle

\begin{center}\small
Code and research records: \url{https://github.com/TejasNaik24/Kleos-Models}\\
Data pipeline: \url{https://github.com/TejasNaik24/Kleos-Training-Data}\qquad
\kleos{}: \url{https://kleos-cs.vercel.app}
\end{center}

\begin{abstract}
% source: headline numbers from docs/experiments.md (EXP:167, 212-220, 545-553, 748-770) and
% docs/experiments/serving-verification-records.md (SVR:137, 193-195); dataset counts from the
% Kleos-Training-Data README (TDR:348-356) and docs/evaluation.md (EVAL:58-64).
\kleos{} is an AI operating system for computer science students whose language models make judgment calls: ranking competing work, resolving conflicting records and deciding when to decline. This paper reports a pre-registered study of teaching one such decision policy to small open-weight models from synthetic data on free GPUs. All 1,350 examples were generated from 18 scenario families and answered by one of 12 explicit policies, so every answer is correct by construction; the 349 examples in a held-out input format form the benchmark (271 answerable, 78 should-decline, 78 perturbation groups). Mistral-Nemo-Instruct 12B (\hermes{}) and Ministral 3 14B (\logos{}: an Instruct variant, then a Reasoning variant trained on policy-derived reasoning traces) were fine-tuned with QLoRA on 820 examples under nine hypotheses registered before the runs, six of them with decision rules, with every deviation logged. Fine-tuning raised Hermes from 0.4755 to 0.8051 ($+0.3295$, 95\% CI 0.3068 to 0.3521, all seven tasks $p<0.001$), replicating a Ministral-8B pilot (0.4744 to 0.8015). Logos v0.0.1 against Hermes on the answerable items was inconclusive ($-0.0168$, cluster 95\% CI $-0.0546$ to $+0.0177$); Logos v0.0.2 was measurably better ($+0.0409$, cluster 95\% CI $+0.0040$ to $+0.0780$, $p=0.031$), while its should-decline gain ($+0.1020$, CI $-0.0018$ to $+0.2062$) was not significant by groups. Each model is one training run, the v0.0.2 comparison confounds three changes, no model produced the held-out JSON format, and abstention was learned by scenario family rather than from evidence. Both models serve from Hugging Face ZeroGPU Spaces and reproduce the frozen evaluation byte for byte on 9 of 9 and 8 of 9 checks; the code, records and data pipeline are public.

\end{abstract}

\section{Introduction}

\kleos{} is an AI operating system for computer science students (\url{https://kleos-cs.vercel.app}). Beneath its planning, memory and research features sits a layer of small decisions that a language model makes many times a day: which of several competing items deserves attention first, which of two contradictory records to trust, whether a request is specific enough to act on, which items belong in a briefing, and which sit in a workspace the user has not opened. These decisions are frequent and consequential, and a general-purpose model prompted with instructions does not reliably make them the way the product requires (the orchestrated base models score 0.4469 to 0.4755 on the benchmark, Section~\ref{sec:results}).

% Figure 1: system overview. Labels only; counts from docs/experiments.md (Status),
% Kleos-Training-Data README (releases) and docs/evaluation.md (benchmark).
\begin{figure}[t]
\centering
\begin{tikzpicture}[
  node distance=4mm and 6mm,
  box/.style={draw, rounded corners=2pt, align=center, font=\footnotesize, inner sep=4pt, text width=30mm, minimum height=11mm},
  wide/.style={box, text width=38mm},
  arr/.style={-{Latex[length=2mm]}, thick}
]
% Row 1: data
\node[box] (cat) {Scenario catalog\\ 18 scenario families,\\ 12 decision policies};
\node[box, right=of cat] (gen) {Generation\\ 1,350 examples; every\\ answer computed by a policy};
\node[box, right=of gen] (gate) {Promotion gates (14)\\ and sealing\\ (content hash)};
\node[wide, right=of gate] (split) {Group split (sha256, seed 42)\\ train 820 / validation 181 /\\ test 349, JSON format held out};
\draw[arr] (cat) -- (gen);
\draw[arr] (gen) -- (gate);
\draw[arr] (gate) -- (split);
% Row 2: models and evaluation
\node[box, below=10mm of cat] (qlora) {QLoRA fine-tuning\\ NF4, $r=16$, T4 GPUs\\ (Colab, Kaggle)};
\node[box, right=of qlora] (arms) {Arms\\ base + orchestration (arm1)\\ vs fine-tuned (arm2)};
\node[box, right=of arms] (bench) {Benchmark, 349 items\\ 271 answerable / 78\\ should-decline; 78 groups};
\node[wide, right=of bench] (verdict) {Pre-registered verdict\\ paired bootstrap, cluster-\\ resampled by group from H8 on};
\draw[arr] (qlora) -- (arms);
\draw[arr] (arms) -- (bench);
\draw[arr] (bench) -- (verdict);
\draw[arr] (split.south west) ++(8mm,0) |- ++(0,-3mm) -| (qlora.north);
% Row 3: serving
\node[box, below=10mm of arms] (pkg) {Deployment package\\ pinned base revision,\\ adapter SHA-256};
\node[box, right=of pkg] (space) {ZeroGPU Space\\ greedy decoding,\\ reproduction check};
\node[box, right=of space] (app) {\kleos{} application\\ Hermes (fast),\\ Logos (shows its reasoning)};
\draw[arr] (verdict.south) |- ++(0,-3mm) -| (pkg.north);
\draw[arr] (pkg) -- (space);
\draw[arr] (space) -- (app);
\end{tikzpicture}
\caption{The pipeline. The data pipeline (top row) lives in the companion repository and seals each release with a content hash; the training, evaluation and serving code (rows two and three) live in the KLEOS Models repository. The test split is the sealed benchmark; it was never used for selection.}
\label{fig:pipeline}
\end{figure}


\kleos{} therefore defines its decisions as an explicit policy. Each candidate item is scored as the product of a deadline term, an evidence weight and an impact weight; items whose scores fall within a relative margin are declined rather than guessed; an explicit statement by the user outranks an inferred record unless it is stale; items outside the active workspace are flagged rather than acted on; and an ambiguous request is answered with a question. The policy is deterministic, so the correct answer to any scenario is computable. That makes two things possible: a synthetic training set in which every answer is correct by construction, and a benchmark on which a model's judgment is measured against the policy rather than against human raters.

This paper asks whether a small open-weight model can learn that policy from such data under conditions a student can afford: free Colab and Kaggle T4 GPUs, QLoRA adapters rather than full fine-tuning, and no proprietary models. It reports two production models. \hermes{} is a QLoRA adapter on Mistral-Nemo-Instruct 12B. \logos{} is an adapter on Ministral 3 14B, in a first version on the Instruct checkpoint and a second on the Reasoning checkpoint, trained on reasoning traces derived from the same policy. Every hypothesis that was run was written down with its metric and decision rule before the run, every deviation from the protocol is dated and logged, and negative and inconclusive results are reported with the same precision as positive ones. The design follows the pre-registration literature \citep{nosek2018preregistration,vanmiltenburg2021preregistering,bertinetto2021preface} and recent guidance on interval estimates for language-model evaluations \citep{miller2024errorbars}.

% source: EXP:167 (pilot), HER:259-275 (Hermes), EXP:545-553 (H8b), EXP:748-770 (H9).
Three findings follow. First, fine-tuning on 820 synthetic examples moved the benchmark score of Hermes from 0.4755 to 0.8051 (difference $+0.3295$, 95\% CI 0.3068 to 0.3521, every one of the seven tasks $p<0.001$), replicating a pilot on Ministral-8B (0.4744 to 0.8015, seven of seven tasks). Second, moving to the stronger Ministral 3 base did not by itself produce a measurably better model: Logos v0.0.1 differed from Hermes by $-0.0168$ on the answerable items (cluster 95\% CI $-0.0546$ to $+0.0177$), an inconclusive result under the pre-registered rule. Third, Logos v0.0.2, trained on the data release that adds policy-derived reasoning traces, was measurably better than Hermes on the answerable items ($+0.0409$, cluster 95\% CI $+0.0040$ to $+0.0780$, $p=0.031$); its gain on the should-decline items did not reach significance by groups. Each model is a single training run, so these intervals cover evaluation noise only, and the v0.0.2 comparison confounds three changes made together.

The contributions are:
\begin{itemize}
  \item A sealed synthetic dataset and benchmark for a product decision policy: 1,350 examples over 18 scenario families and 12 policies, identified by content hashes and rebuilt byte for byte by the public pipeline, whose 349-item test split is byte-identical across the two releases used here (Section~\ref{sec:data}).
  \item A pre-registered protocol for fine-tuning experiments on free compute: dated hypotheses H1 to H9, six of them with decision rules, a cluster-bootstrap amendment adopted before the second model line was trained, and a log of thirteen deviations, D1 to D13, each with its effect on the result (Section~\ref{sec:protocol}, Appendices~\ref{app:deviations} and \ref{app:artifacts}).
  \item The H1 result and its replication across two base models, with sub-score, consistency and abstention analyses (Sections~\ref{sec:results} and \ref{sec:analysis}).
  \item The H8 and H9 comparisons between model lines, with their confounds declared before the runs (Sections~\ref{sec:results} and \ref{sec:analysis}).
  \item A serving path on free ZeroGPU Spaces with pinned base revisions and adapter hashes, verified by byte-for-byte reproduction of the frozen evaluation outputs: 9 of 9 for Hermes and 8 of 9 for Logos (Section~\ref{sec:serving}).
\end{itemize}

Figure~\ref{fig:pipeline} shows how the data pipeline, the training and evaluation harness and the serving path fit together, and Section~\ref{sec:related} places the work among its neighbours. Throughout, every number is copied from a frozen record in the public repository, and each table names the record it reproduces. Analyses are labelled pre-registered, with their hypothesis identifier, or exploratory. Section~\ref{sec:limitations} collects the limitations and Section~\ref{sec:repro} states what can be reproduced, and by whom.

\section{Related work}\label{sec:related}

\paragraph{Parameter-efficient and quantized fine-tuning.}
LoRA \citep{hu2022lora} freezes a pretrained model and trains low-rank updates to selected weight matrices. QLoRA \citep{dettmers2023qlora} stores the frozen weights in a 4-bit NormalFloat format with double quantization and keeps optimizer state in paged memory, which is what lets a 12B or 14B model train on a single T4 (nominally 16\,GB). The 8-bit optimizer \citep{dettmers2022optimizers} and the int8 matrix-multiplication kernels \citep{dettmers2022int8} are the bitsandbytes primitives beneath it; PEFT \citep{mangrulkar2022peft} and Transformers \citep{wolf2020transformers} implement the recipe. This work uses that recipe as published, with no hyperparameter search, and contributes measurements of it on a narrow policy task rather than a new method.

\paragraph{Synthetic and distilled supervision.}
Self-Instruct \citep{wang2023selfinstruct} and Alpaca \citep{taori2023alpaca} generate instruction data from a language model and fine-tune a smaller one on it. Chain-of-thought prompting \citep{wei2022cot} showed that intermediate reasoning text improves answers, and later work trains small models on rationales produced by a larger teacher \citep{ho2023teachers,hsieh2023distilling}; DeepSeek-R1 \citep{guo2025deepseekr1} trains reasoning with reinforcement learning and distils the resulting traces into smaller open models. The data here differs in one respect: both the answers and the reasoning traces are computed by a deterministic policy, not sampled from a teacher model. The supervision therefore carries no model noise and no license entanglement, at the cost that the benchmark measures how well a model learns those policies, not general reasoning.

\paragraph{Statistics for model comparison.}
\citet{dror2018hitchhiker} survey significance testing in natural language processing, and \citet{efron1993bootstrap} is the standard reference for the bootstrap used here. When evaluation items come in groups, resampling must respect the grouping \citep{field2007clustered}, and inference with few clusters over-rejects \citep{cameron2008bootstrap}. \citet{miller2024errorbars} makes the same points for language-model evaluations and recommends clustered standard errors for grouped questions and paired differences when comparing models. \citet{dodge2020finetuning} measured how much fine-tuning results vary with the random seed, and \citet{agarwal2021precipice} argue for interval estimates over point estimates when runs are few. This study adopts paired resampling and, from its second model line on, cluster resampling by benchmark group; its single-run design is its main statistical limitation (Section~\ref{sec:limitations}).

\paragraph{Pre-registration and reporting.}
\citet{nosek2018preregistration} set out the case for registering hypotheses and analysis plans before data are seen; \citet{vanmiltenburg2021preregistering} adapt it to natural language processing and \citet{bertinetto2021preface} to machine learning more broadly. Model cards \citep{mitchell2019modelcards} ask that a released model carry its evaluation conditions and limitations, and \citet{lipton2019troubling} ask that explanation be kept apart from speculation. \citet{altman1995absence} is the reminder that a null result is an estimate with an interval, not evidence of equivalence. The protocol used here is a lightweight version of those proposals: hypotheses, metrics and decision rules committed to a public repository before each run, with the commit identifier as the timestamp, and every outcome reported, including the inconclusive one.

\section{The decision task and the data}\label{sec:data}

\subsection{The tasks and the example format}
% source: task definitions DCM:86-92; system message at index 0, DC:161; domains and formats RP:63-66.
\kleos{} routes its judgment calls through seven decision tasks: notification prioritization (``rank competing items by urgency, evidence and impact''), tool routing (``choose the right tool, or none''), mission-control briefing (``synthesize a structured briefing from context''), context prioritization (``select and order the most relevant context''), recommendation generation (``recommend actions grounded in supplied evidence''), memory-conflict resolution (``resolve contradictions using recency and reliability'') and workspace reasoning (``reason within workspace scoping rules''). An example is a conversation: a system message at index 0 that states the model's role and the rule it must apply; a user message that presents the items in one of four formats, bullets, prose, JSON or a Slack-style thread; and the reference assistant answer. The answer either commits to a ranking, with a one-line rationale per item, or declines and says what it needs; the reference carries the label of the factor that decided it, and in release v0.0.7 every training and validation answer also names that factor in a closing line. The 1,350 examples in the releases used here span four domains (career 335, coursework 341, projects 340, research 334) and the four formats (bullets 473, prose 464, JSON 349, Slack-style 64). Appendix~\ref{app:format} shows the format on a public development fixture.

\subsection{The policy}
% source: formula and constants SC:182-192, 209-215; weights SC:239-252; difficulty SC:269-271; policies SC:196-207.
Each item is scored as
\begin{equation}
\text{score} = \text{deadline\_score} \times \text{evidence\_weight} \times \text{impact\_weight}, \qquad \text{deadline\_score} = \frac{1}{1 + \text{days to the deadline}},
\end{equation}
with evidence weights \code{confirmed} 1.0, \code{corroborated} 0.8, \code{reported} 0.55, \code{single\_source} 0.35, \code{unverified} 0.2 and \code{contradicted} 0.1, and impact weights \code{high} 1.0, \code{medium} 0.6 and \code{low} 0.3. Ties are broken by the item key, so every ranking is total and reproducible. Twelve named policies combine the score with checks; Table~\ref{tab:policies} in Appendix~\ref{app:policies} lists each with what it orders by, when it declines and the deciding factor it names. Three named constants govern the checks: a close call is a relative separation $|a-b|/\max(a,b)$ below 0.15 between the top two scores; an explicit statement is stale after 30 days; and a record that challenges an explicit statement must carry evidence of at least 0.8 (\code{corroborated}). ``Adequate'' evidence means a weight of at least 0.55 (\code{reported}), and an ``explicit statement'' is an item graded \code{confirmed}. An example's difficulty is computed from the policy's own ordering rather than copied from an axis. It depends on the first sort component on which the top two candidates differ: a gating criterion such as scope or adequacy makes the example easy; the primary ranking component makes it easy at a relative gap of at least 0.20 and medium below; a later component makes it medium or hard by the same threshold; and an order that falls to the item key makes it hard. The 0.20 threshold is the smallest relative gap between adjacent evidence grades, $(1.0-0.8)/1.0$. The policy table names eleven distinct deciding-factor labels, of which the benchmark's label set uses nine: six commit labels (\code{deadline}, \code{evidence}, \code{impact}, \code{scope}, \code{explicit\_statement}, \code{least\_privilege}) and five decline labels (\code{insufficient\_separation}, \code{missing\_input}, \code{request\_ambiguous}, \code{stale\_explicit\_conflict}, \code{ask\_before\_crossing}). Appendix~\ref{app:policies} lists which policy produces each label; no record glosses the labels in prose, so that table is assembled from the policy table.

\subsection{Scenario families, perturbations and groups}
% source: families SC:45-71; axes SC:114-120; family size SC:162; perturbation kinds SC:277-289; comparable form SC:217-219; groups RP:63, LC:45-46.
Table~\ref{tab:families} lists the 18 scenario families. Each binds a task to one policy, draws fictional entities from a pool, and varies urgency (deadline pools per urgency class: critical 0, 1 or 2 days; high 1, 2, 3 or 5; medium 4, 7, 10 or 14; low 14, 21, 30 or 45), context length (0, 2 or 4 distractor items) and difficulty spacing (2, 1 or 0 days between adjacent deadlines). A family produces $n_{\text{base}} \times (1 + \sum \text{counts})$ examples: every base example is re-rendered under perturbations of six kinds, \code{paraphrase} (one of three alternative question wordings), \code{evidence\_order} and \code{context\_order} (reordered items), \code{irrelevant\_context} (up to two unused distractor sentences), \code{length} (one distractor added or removed) and \code{schema} (the prompt re-rendered as a Slack-style thread). Every perturbation must keep the decision's comparable form, the triple (ranking, deciding factor, abstained), equal to its base and must differ in content from every other example; generation fails otherwise. A family that reserves the JSON format cannot declare a \code{schema} perturbation, so no group straddles the format holdout. A base example and its perturbations share a group identifier, and the 1,350 examples form 294 such groups. Consistency across a group is therefore a property the data can test, and the group is the unit of dependence that the evaluation statistics must respect (Section~\ref{sec:protocol}).

% source: Kleos-Training-Data docs/scenarios.md, family table (SC:45-65) and the note on JSON-free families (SC:69-71).
\begin{table}[t]
\centering\scriptsize
\caption{The 18 scenario families of releases v0.0.6 and v0.0.7, with the examples each contributes to the release and to the test split, and the out-of-distribution shift it reserves. Task codes: npr notification prioritization, trt tool routing, mcb mission-control briefing, ctx context prioritization, rec recommendation generation, mcr memory-conflict resolution, wsr workspace reasoning. Three families have no JSON-format examples and so contribute nothing to the test split under the format holdout. Source: Kleos-Training-Data, \code{docs/scenarios.md}.}
\label{tab:families}
\begin{tabular}{lllrrl}
\toprule
Family & Task & Policy & Examples & Test & Reserved shift \\
\midrule
\code{ctx.budget\_pressure} & ctx & \code{rank\_by\_relevance\_over\_recency} & 120 & 0 & unseen entities \\
\code{ctx.relevance\_over\_recency} & ctx & \code{rank\_by\_relevance\_over\_recency} & 24 & 8 & reordered evidence \\
\code{mem.explicit\_over\_inferred} & mcr & \code{defer\_to\_explicit\_statement} & 100 & 35 & unseen formats \\
\code{mem.reliability\_over\_recency} & mcr & \code{rank\_by\_reliability\_over\_recency} & 88 & 32 & conflicting evidence \\
\code{mem.stale\_explicit\_conflict} & mcr & \code{defer\_to\_explicit\_statement} & 64 & 20 & conflicting evidence \\
\code{mem.unresolvable\_conflict} & mcr & \code{resolve\_or\_abstain\_on\_support} & 72 & 24 & conflicting evidence \\
\code{brief.across\_workspaces} & mcb & \code{respect\_workspace\_scope} & 100 & 35 & unseen formats \\
\code{brief.what\_needs\_a\_decision} & mcb & \code{rank\_by\_deadline\_then\_evidence} & 24 & 8 & context length shift \\
\code{notif.deadline\_vs\_evidence} & npr & \code{rank\_by\_deadline\_then\_evidence} & 24 & 8 & unseen entities \\
\code{notif.tied\_signals} & npr & \code{rank\_or\_abstain\_when\_close} & 100 & 30 & unseen formats \\
\code{rec.abstain\_without\_evidence} & rec & \code{rank\_or\_abstain\_when\_close} & 18 & 6 & unseen formats \\
\code{rec.verify\_before\_recommending} & rec & \code{verify\_when\_evidence\_weak} & 100 & 35 & unseen formats \\
\code{route.ask\_when\_underspecified} & trt & \code{ask\_when\_request\_ambiguous} & 90 & 30 & unseen formats \\
\code{route.least\_privilege} & trt & \code{prefer\_least\_privilege\_source} & 100 & 0 & unseen entities \\
\code{route.source\_matches\_question} & trt & \code{select\_by\_evidence\_need} & 88 & 32 & unseen formats \\
\code{wsp.absent\_in\_active} & wsr & \code{ask\_before\_crossing\_workspace} & 100 & 0 & unseen entities \\
\code{wsp.cross\_workspace\_request} & wsr & \code{respect\_workspace\_scope} & 120 & 40 & unseen formats \\
\code{wsp.scope\_boundary} & wsr & \code{respect\_workspace\_scope} & 18 & 6 & unseen formats \\
\midrule
Total & & & 1,350 & 349 & \\
\bottomrule
\end{tabular}
\end{table}

\subsection{Splitting and the held-out format}
% source: LC:45-72; SC:69-71, 295-297; configs/datasets/kleos_policy_v006.yaml:10-12, 40-42; SPEC:61-62; STG:110; EVAL:58-64; RM:157 (D2).
A group is never split across train, validation and test. Assignment is a stable hash: each group's rank is the first eight bytes of the SHA-256 of the split seed (42) and the group key, read as a fraction of $2^{64}$. Every release holds out the JSON input format: all 349 JSON-format examples go to the test split, and validation is drawn from the remaining formats, so early stopping never sees the out-of-distribution data the test split is meant to measure. The resulting split is 820 / 181 / 349 (recorded fractions 0.607 / 0.134 / 0.259). Three families (\code{ctx.budget\_pressure}, \code{route.least\_privilege} and \code{wsp.absent\_in\_active}) have no JSON-format examples and contribute nothing to the test split, so 15 families appear there, in 78 groups. A near-duplicate check at a Jaccard similarity of 0.85 or more reports 0 train--test near-duplicates for v0.0.6. The consequence for what the benchmark measures is recorded as deviation D2: test scenarios come from the same families as training, so the benchmark measures transfer across input format, not generalization to new scenarios.

\subsection{Releases and the reasoning traces}
% source: releases TDR:348-362; v0.0.6 changes TCL:154-172; abstain share RP:133-135; v0.0.7 V7:62-63, 93-132, 166-170 and TCL:113-122; test and benchmark hashes SPEC:54, EXP:477.
Releases are immutable and content-hashed. The two used here are v0.0.6 (2026-09-04; content hash \code{3cc9a744\ldots}) and v0.0.7 (2026-10-01; \code{b53afa42\ldots}), both of 1,350 examples in the same 820 / 181 / 349 split and both built from the same catalog with split seed 42. v0.0.6 introduced the relative close-call margin (replacing an absolute margin of 0.08), split the former \code{ask\_when\_underdetermined} policy into \code{verify\_when\_evidence\_weak} and \code{ask\_when\_request\_ambiguous}, and added the stale-statement rule; 338 of its 1,350 targets (25\%) abstain or ask. v0.0.7 keeps every example and decision of v0.0.6, and the identifiers of its JSON-format (test) examples, and adds two things to the training and validation answers only: a reasoning trace (schema 1.1) and a closing line ``What decided it: $\langle$label$\rangle$.'' on every answer, including declined ones. Each trace follows five steps and is capped at 1,200 characters: the rule the family's policy weighs; one line per item, in prompt order, mapping the prompt's wording to the policy's values (the design document's example reads ``Redpine: due in 2 days (0.33), reported (0.55), high (1.0) $\rightarrow$ 0.18''); each check the policy applied and its outcome; why the deciding factor decided it; and the conclusion, either the ranking or the failed check and the question to ask. All 1,001 training and validation answers carry one; the median trace is 560 characters long and the longest 878. The four decline labels that the test split expects thereby enter the training targets for the first time: \code{request\_ambiguous} (45 train / 15 validation), \code{stale\_explicit\_conflict} (32 / 12), \code{missing\_input} (25 / 5) and \code{insufficient\_separation} (22 / 4); the fifth, \code{ask\_before\_crossing}, enters in 100 training and validation answers combined, and no test item expects it. The test split of v0.0.7 is byte-identical to v0.0.6's (SHA-256 \code{a4decaaf\ldots}), and the benchmark derived from it (SHA-256 \code{a11ffad7\ldots}) is the same object for every run in this paper. Earlier releases, v0.0.1 to v0.0.5, record the pipeline's maturation; v0.0.5 removed 32 examples after a human review, described next.

\subsection{Quality and privacy record}
% source: detection layers TPE:20-25; gates STG:83-98, 144 and TDR:84-91; private-fact thresholds TPE:172-178; 32 checks LC:150; approvals PRV:91-94 and TDR:153-154; v0.0.5 review TCL:196-199; labelling quirk V7:67-69; synthetic PRV:42-44, 50.
The pipeline promotes examples through 14 gates, 13 of them mandatory (staging, schema and identifier integrity; four privacy gates; review presence and approval; provenance; corpus deduplication; evaluation leakage; and contract rendering; only the coverage gate can be bypassed). It applies 28 detection rules in three layers: 15 secret patterns that block, 11 personal-data patterns that redact and 2 structural heuristics that flag for review, beside a vault of registered entity literals. A private-fact classifier fails any example scoring 0.75 or more, routes scores from 0.4 to 0.75 to review and passes scores below 0.4. A sealed three-split release passes 32 verification checks. Two limits of this record matter for the paper. First, the approvals in every release to date were filled from the machine review's results; the only completed human review is v0.0.5's sample of 127 examples, of which 32 were rejected (explicit statements not honored, 3; memory conflicts decided on impact, 2; deadlines overridden, 4; ranking where the policy should have asked, 10; unjustified insufficient-separation abstentions, 13). Those categories drove the v0.0.6 policy changes, and no human review of v0.0.6 or v0.0.7 has been completed. Second, in 48 examples, 17 of them in the test split, the deciding factor is the largest raw gap between the top two items even when that gap favors the runner-up, a labelling quirk documented in the v0.0.7 design. No example in any release comes from a real user: every example was generated from the synthetic scenario catalog with fictional entities and captured through an offline mock backend, and the production-observation lane is rejected by the provenance gate.

\section{Models and training}\label{sec:models}

\subsection{Base models}
% source: EXP:90, 473, 647, 939-945; LMD:35-39, 47-53, 353-355; configs/models/*.yaml; RM:334; KMC/models/ministral3_14b_reasoning.yaml:33-37.
Table~\ref{tab:bases} lists the four base checkpoints, all released by Mistral AI \citep{mistral2024nemo,mistral2025mistral3,liu2026ministral3} and all loaded through Transformers 5.16.1 and PEFT 0.20.0. The pilot run used Ministral-8B-Instruct-2410 \citep{mistral2024ministral8b}, released under the Mistral AI Research License, which confines that adapter to research; the run requested revision \code{main}, and the commit it resolved to is not recoverable from the preserved artifacts (deviation D7). \hermes{} is built on Mistral-Nemo-Instruct-2407 \citep{mistral2024nemoinstruct}: 12.2B parameters in 40 layers with hidden size 5,120, MLP width 14,336 and a vocabulary of 131,072, under Apache-2.0, pinned to revision \code{04d8a905\ldots}. \logos{} is built on Ministral 3 14B, whose checkpoints ship a 13,506,073,600-parameter text model beside a 438,958,080-parameter vision encoder and projector. Both Logos runs load the text tower only (40 layers, hidden size 5,120, MLP width 16,384, 32 query and 8 key-value heads of dimension 128, an untied output head) and leave the vision weights on disk. Logos v0.0.1 used the BF16 export of the Instruct checkpoint \citep{mistral2025ministral14binstruct}, revision \code{3cea74c1\ldots}, and v0.0.2 the Reasoning checkpoint \citep{mistral2025ministral14breasoning}, revision \code{51f9210f\ldots}, whose chat template renders an assistant message's \code{reasoning} field as a \code{[THINK]}\ldots\code{[/THINK]} span (token ids 34 and 35, both special) before the answer. The choice of Ministral 3 as the stronger base for H8 rested on published figures from different harnesses and years (MMLU 79.4 against Nemo's 68.0), which the record marks as indicative only.

\begin{table}[t]
\centering\footnotesize
\caption{Base checkpoints. Parameter counts and revisions are as recorded in the model configurations and the experiment record. Both Logos runs load the text tower of the Ministral 3 checkpoint only. Source: \code{docs/experiments.md} and \code{configs/models/}.}
\label{tab:bases}
\begin{tabular}{L{25mm}L{52mm}L{21mm}L{30mm}L{26mm}}
\toprule
Adapter & Checkpoint (Hugging Face) & Revision & Parameters & License \\
\midrule
Ministral-8B pilot (research only) & \code{mistralai/}\allowbreak\code{Ministral-8B-Instruct-2410} & \code{main} (not recoverable, D7) & 8.0B; 36 layers & Mistral AI Research License \\
\hermes{} v0.0.6 & \code{mistralai/}\allowbreak\code{Mistral-Nemo-Instruct-2407} & \code{04d8a905\ldots} & 12.2B; 40 layers & Apache-2.0 \\
\logos{} v0.0.1 & \code{mistralai/}\allowbreak\code{Ministral-3-14B-}\allowbreak\code{Instruct-2512-BF16}, text tower & \code{3cea74c1\ldots} & 13,506,073,600 text (+ 438,958,080 vision, not loaded) & Apache-2.0 \\
\logos{} v0.0.2 & \code{mistralai/}\allowbreak\code{Ministral-3-14B-}\allowbreak\code{Reasoning-2512}, text tower & \code{51f9210f\ldots} & 13,506,073,600 text & Apache-2.0 \\
\bottomrule
\end{tabular}
\end{table}

\subsection{The recipe}
% source: TRN:43-62; AUD:47-53, 72; HER:159-161, 230-232, 498; L1:269-271; EXP:518-523, 697-716; HMD:43; TRN:117-123 (D5); KMC/models/ministral3_14b_reasoning.yaml:46-56.
One QLoRA recipe \citep{dettmers2023qlora,hu2022lora} trained all four adapters (Table~\ref{tab:recipe}). The base model is frozen and stored in 4-bit NormalFloat with double quantization; the compute dtype resolves to float16 on a T4, so every run trained in fp16. LoRA adapters of rank 16 with $\alpha = 32$ and dropout 0.05 attach to the seven projection matrices of every transformer block (\code{q\_proj}, \code{k\_proj}, \code{v\_proj}, \code{o\_proj}, \code{gate\_proj}, \code{up\_proj}, \code{down\_proj}); the output head and the embeddings are excluded. That gives 252 adapted modules and 43,646,976 trainable parameters on the pilot, and 280 modules on the 40-layer models: 57,016,320 parameters for Hermes (per layer $147{,}456 + 98{,}304 + 98{,}304 + 147{,}456 + 3 \times 311{,}296 = 1{,}425{,}408$) and 60,948,480 for both Logos adapters (the wider MLP contributes $3 \times 344{,}064$ per layer). The optimizer is paged 8-bit AdamW \citep{dettmers2022optimizers} at a learning rate of $2 \times 10^{-4}$ with a cosine schedule, a warmup ratio of 0.03 (10 steps), no weight decay and gradient clipping at 1.0. One example per device with 8 accumulation steps gives an effective batch of 8; three epochs over 820 examples are 309 steps. Sequences are capped at 1,024 tokens, which truncates nothing: the longest training example is 440 tokens under the Hermes tokenizer, 448 after padding under Logos v0.0.1's, and 736 under v0.0.2's, where the longest reasoning span is 298 tokens. The loss is computed on assistant tokens only; for v0.0.2 the supervised span is the trace and the answer together. Validation runs every 50 steps for the pilot and Hermes and every 25 for Logos, three checkpoints plus the best are retained, and the adapter with the lowest validation loss is the one shipped. The record describes these settings as engineering defaults, not tuned hyperparameters; no search was run. Appendix~\ref{app:hyper} lists the remaining settings and the software versions of each run.

Two details of the chat templates affected training. Mistral's chat template injects the system message into the last user turn; a training sequence ends on the assistant turn, so the system prompt would have been dropped from every training example. The formatter therefore folds it into the first user turn when the template drops it (deviation D5, fixed before the first run). Ministral 3's tokenizer carries a pre-tokenizer regex that Transformers flags as incorrect; the corrected regex is enabled for both Logos runs and changes the tokenization of 0 of the 1,350 examples. The remaining differences between runs were declared before each run. Logos v0.0.1 differed from Hermes in the base, tokenizer and template; the validation cadence (25 steps instead of 50); strict configuration checking; the allocator setting \code{expandable\_segments}; protection of the best checkpoint from pruning; and a gradient check and a memory probe at the configured batch size before the trainer re-seeds. Logos v0.0.2 differed from v0.0.1 in the base and the data; the generation budget at evaluation (1,024 new tokens instead of 512); the hardware (two T4s, model parallel); adapter selection on a loss that covers the trace as well as the answer; and the results schema.

\begin{table}[t]
\centering\small
\caption{The training recipe shared by all four runs, with the per-run differences. Source: \code{docs/training.md} and the training configurations under \code{configs/training/}.}
\label{tab:recipe}
\begin{tabular}{L{44mm}L{112mm}}
\toprule
Setting & Value \\
\midrule
Quantization & QLoRA: NF4 with double quantization, base frozen; compute dtype \code{auto}, which is float16 on a T4 \\
LoRA & $r = 16$, $\alpha = 32$, dropout 0.05, no bias, default initialization; targets \code{q,k,v,o,gate,up,down\_proj}; \code{lm\_head} and \code{embed\_tokens} excluded \\
Adapted modules / trainable parameters & pilot 252 / 43,646,976; Hermes 280 / 57,016,320; Logos v0.0.1 and v0.0.2 280 / 60,948,480 \\
Optimizer & paged AdamW 8-bit, learning rate $2 \times 10^{-4}$, cosine schedule, warmup ratio 0.03 (10 steps), weight decay 0, max gradient norm 1.0 \\
Batch and epochs & 1 per device $\times$ 8 accumulation steps; 3 epochs over 820 examples $=$ 309 steps \\
Sequence length & 1,024 tokens; nothing truncated \\
Loss & assistant tokens only; for Logos v0.0.2 the reasoning trace and the answer \\
Validation and checkpoints & every 50 steps (pilot, Hermes) or 25 (Logos); keep 3 checkpoints plus the best; ship the lowest validation loss \\
Precision and memory & fp16 with gradient checkpointing; Logos adds a memory probe and a gradient check before training \\
Seed & 42 (\code{deterministic=True}) \\
Software & transformers 5.16.1, peft 0.20.0 and bitsandbytes 0.50.2 in every run; torch 2.11.0+cu128 on Colab; accelerate 1.14.0 recorded for Hermes and both Logos runs (Appendix~\ref{app:hyper}) \\
\bottomrule
\end{tabular}
\end{table}

\subsection{Compute, runs and memory}
% source: hardware HER:69-70, L1:74-75, EXP:652, L2:111-112; software HER:67-68, L1:70-72, L2:72; memory AUD:86, HER:167-172, L1:212-217, L2:102-115; estimator HER:170, LMD:148-149, L1:215; runs AUD:75, HER:84-122, L1:92-115, 131, L2:82-90; GB convention XR:54-55.
The pilot trained on a Colab Tesla T4 (compute capability 7.5, fp16); Hermes and Logos v0.0.1 trained and evaluated on one Colab Tesla T4 (14.56\,GB, compute capability 7.5, CUDA 12.8); and Logos v0.0.2 trained on Kaggle's two T4s with the model's layers spread across both by \code{device\_map: auto} (GPU 0 held the embeddings and layers 0 to 8; GPU 1 held layers 9 to 39, the final norm and the output head). The Colab runs used torch 2.11.0+cu128; Kaggle's torch version is recorded only in the run's private environment file. In these records ``GB'' means $2^{30}$ bytes. Measured peak memory was 9.67\,GB for the pilot, 13.09\,GB for Hermes, 13.70\,GB for Logos v0.0.1 and 10.21\,GB over the whole run for Logos v0.0.2; the memory probe before that run's first step peaked at 10.16\,GB on GPU 1 and 4.75\,GB on GPU 0, against estimates of 8.53 and 6.38\,GB: the total matched, the split did not (finding L-F6). The first memory estimator was about 5\,GB too optimistic for Hermes (7.96 to 8.12\,GB predicted against 13.09\,GB measured, finding H-F10); the rebuilt estimator reproduces the pilot and Hermes peaks within 0.1\% and the Logos v0.0.1 probe within 1.3\%. Findings of this kind, 25 in all across the three registers, are listed in Appendix~\ref{app:findings}.

Table~\ref{tab:runs} summarizes the runs. Hermes trained in two sessions: the Colab runtime died after step 250 while writing a checkpoint, and training resumed eighteen hours later on a different T4 instance from the last complete checkpoint (deviation D8). Logos v0.0.1 needed four sessions and two resumes within the free tier's usage limits (deviation D9). Logos v0.0.2 trained in one Kaggle session. Trainer time for Hermes was 11,222\,s over the two sessions plus 174\,s for the final validation pass, 3\,h\,09\,m\,56\,s in all, at 25.8 to 27.5\,s per step, with nine validation passes totalling 1,645\,s (about 187\,s each in the first session and 174\,s in the second). Logos v0.0.1's record keeps per-session timings but no single total. Logos v0.0.2's trainer reported 15,951\,s (4.43\,h), about 38.5\,s per optimizer step and 51.6\,s per step including validation, in a notebook that ran 17,746\,s with installation, download and a smoke run. The pilot's recorded trainer time is not reconciled with the later runs and is not cited.

\begin{table}[t]
\centering\scriptsize
\caption{The four training runs. ``Val'' is validation loss; for Logos v0.0.2 it covers the trace and the answer. The pilot's \code{train\_runtime} of 1,970.8\,s is in its audit but is not reconciled with the 25.8 to 27.5\,s per step of the later T4 runs, so no duration is cited for it; the Logos v0.0.1 record keeps per-session timings only. Sources: the artifact audit and the three run reports under \code{docs/experiments/}.}
\label{tab:runs}
\begin{tabular}{L{20mm}L{21mm}L{18mm}L{22mm}L{21mm}L{20mm}L{13mm}L{13mm}}
\toprule
Run & Dates (trained) & Hardware & Sessions & Trainer time & Best checkpoint (val) & Final val & Peak memory \\
\midrule
Ministral-8B pilot & 2026-09-09 (frozen 2026-09-15) & 1 $\times$ T4, Colab & not recorded & not reconciled & step 200, epoch 1.95 (0.03987) & 0.04469 & 9.67\,GB \\
\hermes{} v0.0.6 & 2026-09-17 to 09-18 & 1 $\times$ T4, Colab & 2; resumed from step 250 after 18.03\,h & 11,396\,s incl.\ final validation & step 200, epoch 1.95 (0.04343) & 0.04515 & 13.09\,GB \\
\logos{} v0.0.1 & 2026-09-24 to 09-27 & 1 $\times$ T4, Colab & 4; resumed from steps 100 and 250 & per session only; session 4: 2,157.6\,s for 59 steps & step 175, epoch 1.70 (0.03824) & 0.04969 & 13.70\,GB \\
\logos{} v0.0.2 & 2026-10-06 & 2 $\times$ T4, Kaggle, model parallel & 1 & 15,951\,s (4.43\,h) & step 175, epoch 1.70 (0.0293) & 0.0311 & 10.21\,GB (run); 10.16\,GB (GPU 1 probe) \\
\bottomrule
\end{tabular}
\end{table}

% Figure 2: validation loss by step.
% source: Hermes report "Validation curve" (HER:137-145); step 300 was measured once in
%   each session, 0.04511 and 0.04512 (plotted 0.04511).
% source: Logos v0.0.1 report "Validation curve" (L1:172-186).
% source: pilot, docs/experiments.md results log (EXP:873-874): 0.03987 at step 200, 0.04469 at step 309.
% source: Logos v0.0.2 report (L2:88): 0.0293 on the loaded best weights (step 175), 0.0311 at step 309.
\begin{figure}[t]
\centering
\begin{tikzpicture}
\begin{axis}[
  width=0.82\linewidth, height=6.2cm,
  xlabel={Training step (309 steps $=$ 3 epochs)}, ylabel={Validation loss},
  xmin=0, xmax=320, ymin=0.02, ymax=0.17,
  scaled y ticks=false, yticklabel style={/pgf/number format/fixed, /pgf/number format/precision=2},
  ytick={0.03,0.05,0.07,0.09,0.11,0.13,0.15,0.17},
  grid=major, grid style={gray!30},
  legend pos=north east, legend cell align=left,
  legend style={font=\scriptsize, draw=gray!50},
  tick label style={font=\scriptsize}, label style={font=\small},
]
\addplot[color=blue!70!black, mark=*, mark size=1.6pt, thick] coordinates {
  (50,0.09527) (100,0.05465) (150,0.04975) (200,0.04343) (250,0.04437) (300,0.04511) (309,0.04515)};
\addlegendentry{\hermes{} (Mistral-Nemo 12B, v0.0.6)}
\addplot[color=red!75!black, mark=square*, mark size=1.5pt, thick] coordinates {
  (25,0.16277) (50,0.05541) (75,0.06231) (100,0.05292) (125,0.05057) (150,0.03911) (175,0.03824)
  (200,0.04156) (225,0.04605) (250,0.04933) (275,0.04968) (300,0.04973) (309,0.04969)};
\addlegendentry{\logos{} v0.0.1 (Ministral 3 14B Instruct, v0.0.6)}
\addplot[color=gray!60!black, mark=triangle*, only marks, mark size=2.6pt] coordinates {(200,0.03987) (309,0.04469)};
\addlegendentry{Ministral-8B pilot (the two recorded points)}
\addplot[color=black, mark=diamond*, only marks, mark size=2.6pt] coordinates {(175,0.0293) (309,0.0311)};
\addlegendentry{\logos{} v0.0.2 (Reasoning, v0.0.7; trace and answer)}
\addplot[forget plot, color=blue!70!black, mark=o, only marks, mark size=4.5pt, thick] coordinates {(200,0.04343)};
\addplot[forget plot, color=red!75!black, mark=o, only marks, mark size=4.5pt, thick] coordinates {(175,0.03824)};
\end{axis}
\end{tikzpicture}
\caption{Validation loss by training step. Circles mark the checkpoints that were selected on validation loss and shipped: step 200 (epoch 1.95) for Hermes, 0.04343, and step 175 (epoch 1.70) for Logos v0.0.1, 0.03824. Both recorded curves rise through the third epoch, and in the other two runs the final value exceeds the selected one. The pilot and Logos v0.0.2 records keep only the best and final values; v0.0.2's loss covers the reasoning trace as well as the answer. The tokenizers differ between models, so the curves compare in shape, not in level.}
\label{fig:validation}
\end{figure}


\subsection{Checkpoint selection and resume verification}
% source: curves HER:137-154, L1:172-196, EXP:873-874, L2:88, 119-125; loss ratios HER:148-149, L1:193; fp16 skips HER:116, L1:147; gradient check HER:156-161, L1:200; adapters HER:193-196, L1:233-239, EXP:916-917; selection on the trace EXP:710-715; manifests AUD:61-70, HER:57-65.
Figure~\ref{fig:validation} shows the validation curves. In the two runs with a recorded curve, Hermes and Logos v0.0.1, the validation loss reaches its minimum before the end of the second epoch and rises through the third; the pilot and Logos v0.0.2 record only their best and final values, and in both the final value is above the selected one. The pilot and Hermes select step 200 (epoch 1.95), Logos v0.0.1 and v0.0.2 select step 175 (epoch 1.70). The final training loss sits about 10 times below the validation loss for Hermes and about 15 times below it for Logos v0.0.1, which is consistent with a third epoch that memorizes the 820 training targets; selection on validation loss keeps that epoch out of the shipped adapter. Two consequences were declared in advance. Logos' 25-step validation grid could select step 175 where Hermes' 50-step grid could not, and for Logos v0.0.2 the selection criterion covers the trace, which is 58\% of the validation targets by characters, so no answer-only validation loss exists for it. The optimizer skipped 4 steps on fp16 overflow during the Hermes run (through step 305) and 2 during Logos v0.0.1 (both before step 250); the pilot and v0.0.2 records do not report the count.

The interrupted runs were verified rather than assumed faithful. For Hermes, the gradient-check loss before training was bit-identical in both sessions (2.1058197021484375), the scheduler state carried over, and the two independent evaluations of step 300, one per session, agree to $1.4 \times 10^{-5}$ (0.045107 against 0.045121). For Logos v0.0.1, the first resume point was checked by size (368\,MB, a complete checkpoint) and the second file by file (weights, optimizer, scheduler, random-number and scaler state), and the final validation pass on the loaded best weights reproduced the recorded loss exactly. For Hermes and Logos v0.0.1 the exported adapter is byte-identical to the selected checkpoint's weights (SHA-256 \code{dc121fa3\ldots} and \code{f3e8dcdc\ldots}); for the pilot and Logos v0.0.2 the selected checkpoint is established from the recorded validation loss and the retained checkpoints, and v0.0.2's shipped adapter is \code{e49724f6\ldots} (Appendix~\ref{app:artifacts}). Each run's manifest records the hash of its full configuration, the git commit it ran at, the hash of the loaded dataset files and whether any configuration value was adjusted at runtime (none was).

\section{Evaluation protocol}\label{sec:protocol}

\subsection{Benchmark, arms and what was run}
% source: EVAL:33-49, 58-64; RM:63-68; EXP:449-451, 502-504, 564, 670-676; D1 EXP:909.
The benchmark is derived from the sealed test split: 349 items over the seven tasks, in 78 groups, of which 271 are answerable (the reference commits) and 78 should-decline (the reference declines), in 61 and 17 groups respectively. Every item is in the JSON input format, which no training example uses. The protocol defines four arms (Table~\ref{tab:arms}); the pre-registered research question names the prompt-engineered orchestration baseline, so the comparisons in this paper are \code{arm1\_base\_orchestrated} against \code{arm2\_finetuned} (deviation D1: H1 had named the bare \code{arm0\_base}, of which only a 20-example probe exists, scored 0.5659, a biased head-of-file slice that is not comparable with a full run). If that probe were representative, the registered comparison would show a smaller gain than the orchestrated one does; the full bare-base run that would settle this has not been made. The orchestrated arm \code{arm3} was never run, so the interaction hypothesis H7 is only partly measured. Logos v0.0.2 was evaluated on \code{arm2} only: H9 does not ask whether fine-tuning helps, H8a had answered that for the architecture, and a base arm would have cost about 7 GPU hours (estimated) of Kaggle's 30-hour weekly quota. Fine-tuning's own effect on v0.0.2 is therefore not measured.

\begin{table}[t]
\centering\footnotesize
\caption{The evaluation arms and the runs in which each was evaluated. Source: \code{docs/evaluation.md} and the README.}
\label{tab:arms}
\begin{tabular}{L{44mm}L{28mm}L{32mm}L{44mm}}
\toprule
Arm & Model & Prompting & Evaluated in \\
\midrule
\code{arm0\_base} & base & none & no (a 20-example probe only) \\
\code{arm1\_base\_orchestrated} & base & \kleos{} orchestration & pilot, Hermes, Logos v0.0.1 \\
\code{arm2\_finetuned} & base $+$ LoRA adapter & controlled task context & all four runs \\
\code{arm3\_finetuned\_orchestrated} & base $+$ LoRA adapter & \kleos{} orchestration & no \\
\bottomrule
\end{tabular}
\end{table}

\subsection{Decoding and the reasoning split}
% source: configs/evaluation/kleos_policy_v006.yaml:40-52; kleos_logos_v002.yaml:8-23; EVAL:142-160; EXP:727-730; L2:136-145 (L-F7).
Every arm decodes greedily (temperature 0, top-$p$ 1, no sampling, repetition penalty 1) with batch size 1 and one seed, since greedy decoding makes the seed irrelevant. The generation budget is 512 new tokens, raised to 1,024 for Logos v0.0.2 because the worst case computed from the data, twice the longest training trace plus the longest benchmark answer, is $2 \times 298 + 199 = 795$ tokens. For the Reasoning model the completion is split on token id 35 before decoding and only the answer is graded; a completion whose trace never closes has an empty answer, is graded as such, and is marked as cut. The model card for that checkpoint recommends sampling at temperature 1; greedy decoding was kept so that every arm shares one decoding rule and so that the served model can be checked against the evaluation byte for byte (Section~\ref{sec:serving}). One consequence, finding L-F7, is that the stop reason reads ``stop'' whenever the trace closed even if the answer was then cut by the budget; 6 of the 349 v0.0.2 completions were cut that way, all after a closed trace, and were scored as generated.

\subsection{The grader}
% source: EVAL:108-124, 171-199, 250-268; src/kleos_models/evaluation/graders.py:230, 254; HER:559-569; L1:337; D6 EXP:909.
The composite grader \code{kleos\_policy} scores three components against the reference. The \emph{ranking} component is $0.6 \times \text{nDCG} + 0.4 \times \text{top-1 accuracy}$ against the reference ranking \citep{jarvelin2002ndcg}; nDCG was bounded after finding F1 so that a repeated item cannot score above a perfect ranking. The \emph{deciding factor} is read from the \code{deciding\_factor} key of a JSON answer, from a ``What decided it:'' line, or from the first allowed label the answer mentions, and compared with the reference label. \emph{Confidence} is read from an explicit boolean in a JSON answer or, otherwise, from a list of abstention phrasings taken from the \kleos{} corpus, and compared with whether the reference commits. The score is the mean of the components the reference supports. Whether the response is a JSON object (\code{format\_valid}) is reported as a sub-score and excluded from the score (deviation D6): under the format holdout a format failure must be distinguishable from a judgment failure. Kendall's $\tau$ \citep{kendall1938tau} and nDCG are reported as sub-scores. Three faithfulness heuristics (evidence coverage, citation precision and an unsupported-claim rate) are text-level heuristics, not entailment checks; the fabricated-citation heuristic flags 47.9\% of the gold test answers (finding H-F12) and evidence coverage is vacuous because no benchmark item names decisive evidence identifiers (H-F13), so they appear in this paper only beside those floors. Consistency across a group is measured by the agreement rate (the group reached one decision) and the correct-agreement rate (it agreed and was right). The configured grouping unit for the first runs was the scenario family, under which a perfect model scores only 0.600 because 6 of the 15 test families mix different correct labels (finding H-F11); from H8 on, consistency is also reported by the perturbation group, where a perfect model scores 1.000.

\subsection{Subsets and the unlearnable labels}
% source: D3 EXP:909 row; EXP:502-504, 668-672; LMD:323-324; SPEC:39; HER:363-367.
Deviation D3 was discovered during the analysis of the pilot: 78 of the 349 test items, the 78 should-decline items, expect a deciding-factor label that appears zero times in the v0.0.6 training targets (\code{request\_ambiguous} 30, \code{missing\_input} 25, \code{stale\_explicit\_conflict} 20, \code{insufficient\_separation} 3). The cause is a rendering rule: in v0.0.6, decline labels were printed only in JSON-format answers, and the format holdout put every JSON-format example in the test split, so 0 of the 124 training declines carry a label and 492 of the 820 training answers (60\%) have no ``What decided it:'' line at all. For a model trained on v0.0.6 the deciding-factor component is therefore capped at 0.7765 and the score at 0.9255 on the full benchmark, and the deciding factor is read on the answerable items as well as overall (0.816 on the 271 answerable items for the pilot, 0.6332 overall). This affects both arms of every v0.0.6 comparison identically, so H1 stays fair, but it is the reason the two cross-model hypotheses, H8b and H9, name the 271 answerable items as their primary population, with the 78 should-decline items as the first secondary row. v0.0.7 trains the four decline labels for the first time, and the should-decline row is the one it was built to move; H9 kept the answerable subset as primary on the stated ground that switching after the fact would ``pick the population most likely to flatter the result.''

\subsection{Statistics}
% source: EXP:818, 825-836; EVAL:215-217, 296, 318; metrics.py:330-366, 380, 395-436; EXP:489, 499-500, 665; XR:57-58; EXP:561-562, 765-766.
Differences are paired per item. A difference's 95\% interval comes from 2,000 bootstrap resamples with seed 42, taking the resampled differences at positions $\lfloor 0.025 B \rfloor$ and $\lfloor 0.975 B \rfloor$ of the sorted list; the two-sided $p$-value is twice the fraction of resamples on the far side of zero. When no resample crossed, the comparison tool prints $p = 0.0$ (H1) or $p < 0.0005$ (H8 onward, which is $1/B$); under the doubled one-sided fraction the two-sided bound is $2/B$, so this paper reports every such case as $p < 0.001$. A protocol amendment of 2026-09-24, adopted before any Logos training, changed the resampling unit: from H8 on, every interval resamples the 78 groups rather than the 349 items, because perturbations of one case are not independent and item-level intervals are too narrow (finding H-F14; \citealp{field2007clustered,miller2024errorbars}). Item-level intervals are still printed beside the cluster ones, no earlier result was re-decided, and intervals over fewer than 5 groups are reported as not estimable, which affects the 8-item \code{context\_prioritization} task. With 61 answerable and 17 should-decline groups, cluster inference is itself approximate; with few clusters it tends to over-reject \citep{cameron2008bootstrap}. Per-task rows are reported without correction for multiple comparisons and are not decisive. The decision rules are fixed in advance. For a fine-tuning hypothesis (H1, H8a) a task improves only if its interval excludes zero, and the hypothesis is supported if at least one task improves and none regresses. For a cross-model hypothesis (H8b, H9) the verdict on the paired cluster interval of the primary difference is \emph{better} if the interval lies entirely above zero, \emph{worse} if entirely below, \emph{equivalent} if entirely within $\pm 0.02$, and \emph{inconclusive} otherwise. Every result in this paper is reported as an estimate with its interval; a null is never read as equivalence \citep{altman1995absence}.

\subsection{Pre-registration}
% source: Table 4 from EXP:44-83, 136-146, 236-244, 273-281, 356-362, 484-500, 660-676; dates and commits from the git history (first commit fd87931, 2026-08-14) and the commit map XR:106-132; fixed protocol EXP:812-821.
Table~\ref{tab:hypotheses} lists the nine hypotheses with their metrics, decision rules, registration dates and outcomes. H1 to H7 and the fixed protocol (held-out splits only, greedy decoding identical across arms, a benchmark and graders fixed before the run, per-task reporting with no blended aggregate, recorded seeds, a leakage check before training) were committed on 2026-08-14 with the repository's first commit. H5 was re-scoped and H6 closed on 2026-09-15, after the pilot and before Hermes was trained. H8 and the clustered-interval amendment were committed on 2026-09-24, before any Logos training; deviation D11 was declared on 2026-09-29 before the arm it concerns ran; H9 was committed on 2026-10-06, before any Logos v0.0.2 training. In each case the commit that adds the entry is the registration and each run's manifest records the commit it ran at, so the order of registration and run is checkable from the public history (commit identifiers here are the post-rewrite ones; Appendix~\ref{app:artifacts} maps them to the identifiers the frozen records cite). Thirteen deviations, D1 to D13, are dated and logged with their effect on the result (Appendix~\ref{app:deviations}).

\begin{table}[t]
\centering\scriptsize
\caption{The pre-registered hypotheses. ``Registered'' is the date and post-rewrite commit of the entry; H5 to H7 carry no decision rule in the record. The record reports no interval for H3, although its rule names one. Source: \code{docs/experiments.md}, Status table and hypothesis entries.}
\label{tab:hypotheses}
\begin{tabular}{L{6mm}L{32mm}L{32mm}L{38mm}L{16mm}L{28mm}}
\toprule
Id & Question & Primary metric & Decision rule & Registered & Status \\
\midrule
H1 & Fine-tuning beats the prompt-engineered baseline & per-task grader score on the held-out test split (registered as arm0 vs arm2; run as arm1 vs arm2, D1) & improvement counts only if the paired bootstrap 95\% CI excludes zero; per task, no blended aggregate & 2026-08-14, \code{fd87931} & Supported on the pilot (2026-09-15, with D1 and D2); replicated by Hermes (2026-09-22) \\
H2 & Gains generalize out of distribution & \code{ood\_score} beside \code{in\_distribution\_score}; generalization gap & gains generalize if the OOD delta CI excludes zero & 2026-08-14 & Not measurable: the v0.0.6 benchmark is entirely out of distribution (2026-09-15) \\
H3 & Decisions survive irrelevant perturbations & \code{correct\_agreement\_rate}; flips per perturbation kind & consistency improves if the delta CI excludes zero & 2026-08-14 & Improved, still poor (pilot 2026-09-15; reproduced exactly by Hermes 2026-09-22); no interval computed \\
H4 & No loss of general capability & \code{relative\_delta} on a fixed capability suite & regression if $< -0.01$ with a CI excluding zero & 2026-08-14 & Not run \\
H5 & The effect depends on scale & --- & --- & 2026-08-14; re-scoped 2026-09-15 & Not run; re-scoped to two Mistral models \\
H6 & Transfer across model families & --- & --- & 2026-08-14; closed 2026-09-15 & Closed: not testable with one model family \\
H7 & Fine-tuning and orchestration interact & --- & --- & 2026-08-14 & Partially measured: arm1 and arm2 only, in three runs (2026-10-01) \\
H8a & Fine-tuning helps Logos v0.0.1 & per-task \code{kleos\_policy}, Logos arm1 vs arm2 & a task improves only if its cluster CI excludes zero; supported if at least one improves and none regresses & 2026-09-24, \code{81e3e9a} & Supported (2026-10-01) \\
H8b & Logos v0.0.1 is better than Hermes (primary) & mean \code{kleos\_policy} on the 271 answerable items, paired & better / worse / equivalent ($\pm 0.02$) / inconclusive, by the cluster CI of Logos $-$ Hermes & 2026-09-24, \code{81e3e9a} & Inconclusive (2026-09-29) \\
H9 & Logos v0.0.2, trained to think, is better than Hermes (primary) & as H8b & as H8b & 2026-10-06, \code{6d7e450} & Supported: better, $+0.0409$, cluster 95\% CI $+0.0040$ to $+0.0780$ (2026-10-06) \\
\bottomrule
\end{tabular}
\end{table}

\section{Results}\label{sec:results}

Every result in this section is pre-registered unless marked as a re-report or as descriptive. Scores are means of the \code{kleos\_policy} grader on the 349-item benchmark. Intervals are 95\% bootstrap intervals, item-level for the H1 comparisons and cluster-level by perturbation group from H8 on, as the protocol fixes. Table~\ref{tab:main} collects the main numbers, Table~\ref{tab:per-task} the per-task results, Table~\ref{tab:head-to-head} the comparisons between fine-tuned models and Table~\ref{tab:subscores} the grader sub-scores; Figures~\ref{fig:main} to \ref{fig:forest} draw them.

% Figure 3: mean kleos_policy score on the 349-item benchmark, base + orchestration (arm1)
% vs fine-tuned (arm2), with the recorded example-level 95% bootstrap intervals.
% source: pilot, docs/experiments.md H1 corrected scores (EXP:167): 0.4744 / 0.8015; no corrected
%   interval is recorded for the baseline, so neither pilot bar carries one.
% source: Hermes report (HER:259-262): arm1 0.4755 [0.4532-0.4978], arm2 0.8051 [0.7813-0.8288].
% source: Logos v0.0.1 report (L1:291-294): arm1 0.4469 [0.4219-0.4720], arm2 0.7896 [0.7644-0.8147].
% source: Logos v0.0.2, docs/experiments.md H9 result (EXP:790-791): arm2 0.8596, descriptive, no arm1.
% The error-bar lengths are the distances from each mean to its recorded interval endpoints.
\begin{figure}[t]
\centering
\begin{tikzpicture}
\begin{axis}[
  ybar, bar width=11pt,
  width=0.8\linewidth, height=6cm,
  ylabel={Mean \code{kleos\_policy} score}, ymin=0, ymax=1.0,
  symbolic x coords={Ministral-8B pilot, Hermes, Logos v0.0.1, Logos v0.0.2},
  xtick={Ministral-8B pilot, Hermes, Logos v0.0.1, Logos v0.0.2}, enlarge x limits=0.18,
  ytick={0,0.2,0.4,0.6,0.8,1.0},
  ymajorgrids, grid style={gray!30},
  legend style={at={(0.5,1.03)}, anchor=south, legend columns=2, font=\scriptsize, draw=gray!50}, legend cell align=left,
  tick label style={font=\scriptsize}, label style={font=\small},
  nodes near coords, nodes near coords style={font=\tiny, yshift=5pt, /pgf/number format/.cd, fixed, precision=4},
  error bars/y dir=both, error bars/y explicit,
  error bars/error bar style={black, thin},
]
\addplot[fill=gray!40, draw=gray!60!black] plot[error bars/.cd, y dir=both, y explicit] table[x=model, y=score, y error plus=ep, y error minus=em] {
model score ep em
{Ministral-8B pilot} 0.4744 0 0
{Hermes} 0.4755 0.0223 0.0223
{Logos v0.0.1} 0.4469 0.0251 0.0250
};
\addlegendentry{Base model with orchestration (arm1)}
\addplot[fill=blue!60!black, draw=blue!30!black] plot[error bars/.cd, y dir=both, y explicit] table[x=model, y=score, y error plus=ep, y error minus=em] {
model score ep em
{Ministral-8B pilot} 0.8015 0 0
{Hermes} 0.8051 0.0237 0.0238
{Logos v0.0.1} 0.7896 0.0251 0.0252
{Logos v0.0.2} 0.8596 0 0
};
\addlegendentry{Fine-tuned (arm2)}
\end{axis}
\end{tikzpicture}
\caption{Mean \code{kleos\_policy} score on the 349-item benchmark. Error bars are the recorded item-level 95\% bootstrap intervals (Hermes and Logos v0.0.1 only; the pilot baseline was re-graded after its interval was recorded, and the Logos v0.0.2 overall mean is descriptive). Logos v0.0.2 has no base arm. Item-level intervals treat the 349 items as independent and are narrower than the cluster intervals used from H8 on (Section~\ref{sec:protocol}).}
\label{fig:main}
\end{figure}

% Figure 4: per-task gain (arm2 - arm1) for the three runs with a base arm.
% source: pilot, docs/experiments.md H1 corrected table (EXP:157-165).
% source: Hermes report per-task table (HER:266-275).
% source: Logos v0.0.1, docs/experiments.md H8a table (EXP:594-603).
% Task order is the data-contract order; n is the test count per task.
\begin{figure}[t]
\centering
\begin{tikzpicture}
\begin{axis}[
  ybar, bar width=6.5pt,
  width=\linewidth, height=6cm,
  ylabel={Gain in mean score (fine-tuned $-$ base)}, ymin=0, ymax=0.52,
  symbolic x coords={notif. ($n=38$), tool routing ($n=62$), briefing ($n=43$), context ($n=8$), recommend. ($n=41$), memory ($n=111$), workspace ($n=46$)},
  xtick=data, x tick label style={font=\scriptsize, rotate=20, anchor=east},
  enlarge x limits=0.09,
  ymajorgrids, grid style={gray!30},
  legend pos=north west, legend cell align=left, legend style={font=\scriptsize, draw=gray!50},
  tick label style={font=\scriptsize}, label style={font=\small},
]
\addplot[fill=gray!40, draw=gray!60!black] coordinates {
  (notif. ($n=38$),0.2617) (tool routing ($n=62$),0.3262) (briefing ($n=43$),0.4385) (context ($n=8$),0.3870)
  (recommend. ($n=41$),0.1284) (memory ($n=111$),0.3412) (workspace ($n=46$),0.4108)};
\addlegendentry{Ministral-8B pilot (v0.0.6)}
\addplot[fill=blue!60!black, draw=blue!30!black] coordinates {
  (notif. ($n=38$),0.1904) (tool routing ($n=62$),0.2750) (briefing ($n=43$),0.4035) (context ($n=8$),0.4065)
  (recommend. ($n=41$),0.1971) (memory ($n=111$),0.3798) (workspace ($n=46$),0.4323)};
\addlegendentry{\hermes{} (v0.0.6)}
\addplot[fill=red!65!black, draw=red!35!black] coordinates {
  (notif. ($n=38$),0.3054) (tool routing ($n=62$),0.2479) (briefing ($n=43$),0.4094) (context ($n=8$),0.3792)
  (recommend. ($n=41$),0.0643) (memory ($n=111$),0.4586) (workspace ($n=46$),0.4005)};
\addlegendentry{\logos{} v0.0.1 (v0.0.6)}
\end{axis}
\end{tikzpicture}
\caption{Per-task gain from fine-tuning on the 349-item benchmark, by task (test items per task in parentheses). Every pilot and Hermes gain was significant under the pre-registered item-level bootstrap ($p<0.05$ and $p<0.001$ respectively); for Logos v0.0.1, five tasks improved under the cluster bootstrap, \code{recommendation\_generation} did not reach significance, and \code{context\_prioritization} (8 items, fewer than 5 groups) was not estimable by groups. Per-task intervals are not recorded, so no error bars are drawn (Table~\ref{tab:per-task}).}
\label{fig:tasks}
\end{figure}

% Figure 5: the six recorded paired cluster-bootstrap differences between fine-tuned models.
% source: H8b primary and should-decline rows, docs/experiments.md (EXP:545-549, 564):
%   answerable -0.0168 [-0.0546, +0.0177]; should-decline -0.0111 [-0.0429, +0.0195].
% source: H9 primary row (EXP:748-752) and should-decline row (EXP:768-770; L2:166):
%   answerable +0.0409 [+0.0040, +0.0780]; should-decline +0.1020 [-0.0018, +0.2062].
% source: Logos v0.0.1 -> v0.0.2, Logos v0.0.2 report section 8 (L2:204-210):
%   answerable +0.0577 [+0.0234, +0.0952]; should-decline +0.1131 [+0.0028, +0.2296].
% Answerable rows resample 61 groups, should-decline rows 17. Error-bar lengths are the
% distances from each estimate to its recorded interval endpoints.
\begin{figure}[t]
\centering
\begin{tikzpicture}
\begin{axis}[
  width=0.88\linewidth, height=6.4cm,
  xlabel={Paired difference in mean \code{kleos\_policy} score (95\% cluster-bootstrap interval)},
  xmin=-0.08, xmax=0.25,
  xtick={-0.08,-0.04,0,0.04,0.08,0.12,0.16,0.20,0.24},
  scaled x ticks=false, xticklabel style={/pgf/number format/fixed, /pgf/number format/precision=2},
  ymin=0.4, ymax=6.6, y dir=reverse,
  ytick={1,2,3,4,5,6},
  yticklabels={
    {H8b: Logos v0.0.1 $-$ Hermes, answerable (61 groups)},
    {H8b: Logos v0.0.1 $-$ Hermes, should-decline (17 groups)},
    {H9: Logos v0.0.2 $-$ Hermes, answerable (61 groups)},
    {H9: Logos v0.0.2 $-$ Hermes, should-decline (17 groups)},
    {Logos v0.0.2 $-$ v0.0.1, answerable (61 groups)},
    {Logos v0.0.2 $-$ v0.0.1, should-decline (17 groups)}},
  y tick label style={font=\scriptsize, align=right},
  tick label style={font=\scriptsize}, label style={font=\small},
  xmajorgrids, grid style={gray!30},
  clip=false,
]
\fill[gray!18] (axis cs:-0.02,0.4) rectangle (axis cs:0.02,6.6);
\draw[dashed, gray!60!black] (axis cs:0,0.4) -- (axis cs:0,6.6);
\node[font=\tiny, gray!60!black, anchor=south] at (axis cs:0,0.4) {$\pm 0.02$ equivalence margin};
\addplot[only marks, mark=*, mark size=2.2pt, color=black,
  error bars/.cd, x dir=both, x explicit, error bar style={thick}, error mark options={rotate=90, mark size=2.5pt}]
  table[x=est, y=row, x error plus=ep, x error minus=em] {
est row ep em
-0.0168 1 0.0345 0.0378
-0.0111 2 0.0306 0.0318
0.0409 3 0.0371 0.0369
0.1020 4 0.1042 0.1038
0.0577 5 0.0375 0.0343
0.1131 6 0.1165 0.1103
};
\end{axis}
\end{tikzpicture}
\caption{Paired differences between fine-tuned models, resampled by benchmark group (2,000 resamples). The primary rows are the answerable ones: H8b is inconclusive because its interval crosses zero and is not contained in the $\pm 0.02$ equivalence margin, and H9 is ``better'' because its interval lies entirely above zero. The should-decline rows rest on 17 groups; only the v0.0.1 to v0.0.2 row excludes zero.}
\label{fig:forest}
\end{figure}


% source: pilot EXP:167, 182-198; Hermes HER:259-262, HMD:110-111, L1:416-418, HMD:135-136; Logos v0.0.1 L1:291-294, EXP:605-613; Logos v0.0.2 EXP:748-752, 787-791, L2:154-176.
\begin{table}[t]
\centering\scriptsize
\caption{Main results on the 349-item benchmark. Base and fine-tuned intervals are item-level; the answerable and should-decline columns are re-reports by group with cluster intervals where recorded. The pilot's baseline was re-graded after its interval was recorded, so neither pilot figure carries one; Logos v0.0.2 has no base arm and its overall mean is descriptive. Consistency is the agreement rate by perturbation group, where a perfect model scores 1.000. Sources: \code{docs/experiments.md} (H1, H8, H9 entries) and the run reports.}
\label{tab:main}
\begin{tabular}{L{19mm}L{21mm}L{21mm}L{26mm}L{27mm}L{20mm}L{16mm}}
\toprule
Model & Base, arm1 & Fine-tuned, arm2 & Difference & Answerable (271 items, 61 groups), arm2 & Should-decline (78 items, 17 groups), arm2 & Consistency, arm2 (all / answerable) \\
\midrule
Ministral-8B pilot & 0.4744 & 0.8015 & $+0.3271$; 7 of 7 tasks $p<0.05$ & not reported & not reported & not reported \\
\hermes{} & 0.4755 \ci{0.4532}{0.4978} & 0.8051 \ci{0.7813}{0.8288} & $+0.3295$ \ci{0.3068}{0.3521}, $p<0.001$ & 0.8976 \ci{0.8634}{0.9301} (base 0.5570) & 0.4835 (base 0.1923) & 0.769 / 0.770 (base 0.346) \\
\logos{} v0.0.1 & 0.4469 \ci{0.4219}{0.4720} & 0.7896 \ci{0.7644}{0.8147} & $+0.3427$ (descriptive); answerable $+0.3952$ \ci{+0.3443}{+0.4438} & 0.8808 \ci{0.8381}{0.9191} (base 0.4857 \ci{0.4445}{0.5261}) & 0.4724 (base 0.3123; $+0.1601$ \ci{+0.0712}{+0.2409}) & 0.833 / 0.853 (base 0.154 / 0.164) \\
\logos{} v0.0.2 & not run & 0.8596 (descriptive) & --- & 0.9385 \ci{0.9083}{0.9663} & 0.5855 & 0.833 / 0.869 \\
\bottomrule
\end{tabular}
\end{table}

\subsection{H1: the pilot and its replication}
% source: EXP:157-198 (pilot per task, F1 re-grade); HER:259-320 (Hermes arms, per task, sub-scores); HER:384-437 and EXP:103-108 (descriptive comparison); HMD:110-111, L1:416-418, HMD:135-136 (subsets and consistency by group).
The pilot run on Ministral-8B compared the orchestrated base against the fine-tuned adapter on all 349 items. Its first report scored the baseline at 0.5231. An artifact audit then found that nDCG could exceed 1.0 on a ranking that repeats an item (finding F1), and 60\% of the baseline's answers repeated one. Bounding nDCG re-graded 210 of the 349 baseline responses and none of the fine-tuned ones, moved the baseline to 0.4744, left the fine-tuned score at 0.8015, and raised the count of tasks with a significant improvement from six to seven. The corrected gap is $+0.3271$, with every task improved at $p<0.05$ under the pre-registered item-level bootstrap (Table~\ref{tab:per-task}, first block); parse failures fell from 13 to 0 (descriptive).

\hermes{} replicated the result on a different base one week later: 0.4755 \ci{0.4532}{0.4978} to 0.8051 \ci{0.7813}{0.8288}, a paired difference of $+0.3295$ \ci{0.3068}{0.3521}. The comparison tool prints $p = 0.0$, meaning that none of the 2,000 resamples crossed zero, so the correct statement is $p<0.001$. Every task improved with $p<0.001$ (Table~\ref{tab:per-task}, second block; Figure~\ref{fig:tasks}), from $+0.1904$ on notification prioritization to $+0.4323$ on workspace reasoning; the 8-item context-prioritization task moved from 0.5935 to 1.0000 (interval 0.2612 to 0.5352). The sub-scores (Table~\ref{tab:subscores}) show where the gain comes from: the ranking component rose from 0.5985 to 0.8822 (nDCG 0.6174 to 0.9527, top-1 accuracy 0.5702 to 0.7765, Kendall's $\tau$ 0.5449 to 0.8479), the deciding factor from 0.0860 to 0.6132 and confidence from 0.7421 to 0.9198. \code{format\_valid} is 0.0000 in both arms with 0 parse failures: neither arm produced the held-out JSON format, and the scores measure judgment on prose answers. Descriptively, fine-tuning also shortened the answers, from a mean of 226 completion tokens to 94 (medians 217 and 91); one baseline answer hit the 512-token cap and no fine-tuned answer did; mean latency on the T4 fell from 24.7\,s to 16.1\,s. Re-reported by group after the H8 amendment, Hermes scores 0.8976 \ci{0.8634}{0.9301} on the 271 answerable items (base 0.5570) and 0.4835 on the 78 should-decline items (base 0.1923), and its agreement rate by group is 0.769 (base 0.346).

The two fine-tuned models land within 0.004 of each other overall (0.8015 and 0.8051) and the two baselines within 0.002; the per-task differences between the adapters range from $-0.0766$ on notification prioritization to $+0.1076$ on context prioritization. The record treats these as leads, not results: scale, architecture and pretraining are confounded between the two bases, and the agreement is cross-base consistency, not a replication of one configuration across seeds.

% source: pilot EXP:157-165; Hermes HER:266-275; Logos v0.0.1 EXP:594-603.
\begin{table}[t]
\centering\scriptsize
\caption{Per-task results for the three runs with a base arm. Verdicts for the pilot and Hermes are under the pre-registered item-level bootstrap; verdicts for Logos v0.0.1 are under the cluster bootstrap by group (the item-level verdict is given where it differs). Per-task intervals are not recorded except the one shown. Sources: \code{docs/experiments.md} H1 and H8a tables; Hermes run report.}
\label{tab:per-task}
\begin{tabular}{L{40mm}R{8mm}R{14mm}R{16mm}R{14mm}L{54mm}}
\toprule
Task & $n$ & Base & Fine-tuned & Difference & Verdict \\
\midrule
\multicolumn{6}{l}{\emph{Ministral-8B pilot, corrected after F1 (2026-09-15)}} \\
mission-control briefing & 43 & 0.5323 & 0.9708 & $+0.4385$ & improved, $p<0.05$ \\
workspace reasoning & 46 & 0.5228 & 0.9336 & $+0.4108$ & improved, $p<0.05$ \\
context prioritization & 8 & 0.5054 & 0.8924 & $+0.3870$ & improved, $p<0.05$ \\
memory-conflict resolution & 111 & 0.4781 & 0.8192 & $+0.3412$ & improved, $p<0.05$ \\
tool routing & 62 & 0.3919 & 0.7181 & $+0.3262$ & improved, $p<0.05$ \\
notification prioritization & 38 & 0.6253 & 0.8871 & $+0.2617$ & improved, $p<0.05$ \\
recommendation generation & 41 & 0.3285 & 0.4569 & $+0.1284$ & improved, $p<0.05$ \\
overall & 349 & 0.4744 & 0.8015 & $+0.3271$ & 7 of 7 tasks improved \\
\midrule
\multicolumn{6}{l}{\emph{\hermes{} (2026-09-22)}} \\
workspace reasoning & 46 & 0.5190 & 0.9512 & $+0.4323$ & improved, $p<0.001$ \\
context prioritization & 8 & 0.5935 & 1.0000 & $+0.4065$ & improved, $p<0.001$; interval 0.2612 to 0.5352 \\
mission-control briefing & 43 & 0.5681 & 0.9716 & $+0.4035$ & improved, $p<0.001$ \\
memory-conflict resolution & 111 & 0.4756 & 0.8553 & $+0.3798$ & improved, $p<0.001$ \\
tool routing & 62 & 0.3929 & 0.6679 & $+0.2750$ & improved, $p<0.001$ \\
recommendation generation & 41 & 0.2974 & 0.4946 & $+0.1971$ & improved, $p<0.001$ \\
notification prioritization & 38 & 0.6201 & 0.8105 & $+0.1904$ & improved, $p<0.001$ \\
overall & 349 & 0.4755 & 0.8051 & $+0.3295$ & interval 0.3068 to 0.3521, $p<0.001$ \\
\midrule
\multicolumn{6}{l}{\emph{\logos{} v0.0.1, H8a (2026-10-01), verdicts by group}} \\
memory-conflict resolution & 111 & 0.3763 & 0.8350 & $+0.4586$ & improved, $p<0.001$ \\
mission-control briefing & 43 & 0.5552 & 0.9646 & $+0.4094$ & improved, $p<0.001$ \\
workspace reasoning & 46 & 0.5366 & 0.9370 & $+0.4005$ & improved, $p<0.001$ \\
context prioritization & 8 & 0.6204 & 0.9995 & $+0.3792$ & not estimable by group (8 items); improved by item, $p<0.001$ \\
notification prioritization & 38 & 0.5888 & 0.8942 & $+0.3054$ & improved, $p<0.001$ \\
tool routing & 62 & 0.4021 & 0.6501 & $+0.2479$ & improved, $p<0.001$ \\
recommendation generation & 41 & 0.3263 & 0.3906 & $+0.0643$ & not significant \\
overall & 349 & 0.4469 & 0.7896 & $+0.3427$ & descriptive; 5 tasks improved, none regressed: supported \\
\bottomrule
\end{tabular}
\end{table}

\subsection{H8: a stronger base}
% source: H8a EXP:594-613, L1:291-300, 321-327; H8b EXP:545-576, L1:395-437.
\paragraph{H8a (supported).} Fine-tuning helped the Ministral 3 base as it had helped the others: Logos v0.0.1 moved from 0.4469 \ci{0.4219}{0.4720} to 0.7896 \ci{0.7644}{0.8147} overall. Under the cluster bootstrap five tasks improved with $p<0.001$ and none regressed (Table~\ref{tab:per-task}, third block); recommendation generation ($+0.0643$) was not significant, and the 8-item context-prioritization task was not estimable by group. On the answerable subset the paired cluster difference is $+0.3952$ \ci{+0.3443}{+0.4438} (0.4857 to 0.8808) and on the should-decline subset $+0.1601$ \ci{+0.0712}{+0.2409} (0.3123 to 0.4724); agreement by group rose from 0.154 to 0.833. One caveat is recorded beside these gains: 41 of the base's 349 answers (11.7\%) were cut at the 512-token budget, against none of the adapter's (mean completion length 348 against 95 tokens), which may inflate the base-to-adapter differences.

\paragraph{H8b (inconclusive).} On the primary population Hermes scores 0.8976 \ci{0.8634}{0.9301} and Logos v0.0.1 0.8808 \ci{0.8381}{0.9191}; the paired cluster difference over 61 groups is $-0.0168$ \ci{-0.0546}{+0.0177} (Table~\ref{tab:head-to-head}, Figure~\ref{fig:forest}). The interval crosses zero and is not contained in the $\pm 0.02$ equivalence margin, so the pre-registered verdict is inconclusive: a Hermes lead of up to 0.0546 is not ruled out, nor is a Logos lead of up to 0.0177. The secondary rows point the same way. On the should-decline subset the difference is $-0.0111$ \ci{-0.0429}{+0.0195}. Per task, recommendation generation regressed from 0.4946 to 0.3906 ($-0.1040$; $p = 0.006$ by group, $p<0.001$ by item) and workspace reasoning by $-0.0142$ ($p = 0.014$ by group, not significant by item); notification prioritization rose by $+0.0837$ ($p = 0.007$ by item, not significant by group); the other three estimable tasks' group intervals included zero (per-task intervals are not recorded). Agreement by group was higher for Logos (0.833 against 0.769; 0.853 against 0.770 on answerable groups), with no interval computed (descriptive). The overall means, 0.8051 and 0.7896, are descriptive.

\subsection{H9: a Logos trained to think}
% source: EXP:748-791; L2:136-145, 154-176, 185-194, 204-210.
\paragraph{H9 (supported: better).} Logos v0.0.2 scores 0.9385 \ci{0.9083}{0.9663} on the answerable items against Hermes' 0.8976 \ci{0.8634}{0.9301}; the paired cluster difference over 61 groups is $+0.0409$ \ci{+0.0040}{+0.0780}, $p = 0.031$. The interval lies entirely above zero, so the pre-registered verdict is better; the record notes that the lower bound is close to zero. The first secondary row, the one the v0.0.7 data was built to move, did not reach significance by group: on the should-decline items Logos v0.0.2 scores 0.5855 against 0.4835, $+0.1020$ \ci{-0.0018}{+0.2062}, $p = 0.054$ over 17 groups. By task, notification prioritization ($+0.1324$, $p<0.001$), recommendation generation ($+0.0893$, $p = 0.002$) and memory-conflict resolution ($+0.0737$, $p = 0.048$) improved; tool routing ($+0.0330$), workspace reasoning ($+0.0172$) and mission-control briefing ($-0.0153$) had group intervals that included zero; context prioritization ($-0.0002$) was not estimable. Descriptively, agreement by group rose from 0.769 to 0.833 over all groups and from 0.770 to 0.869 over the answerable ones. The overall mean, 0.8596, is descriptive: v0.0.2 has no base arm.

\paragraph{Logos v0.0.1 to v0.0.2 (secondary).} Against its own predecessor, v0.0.2 gains $+0.0577$ \ci{+0.0234}{+0.0952} on the answerable items ($p<0.001$) and $+0.1131$ \ci{+0.0028}{+0.2296} on the should-decline items ($p = 0.046$); recommendation generation ($+0.1933$, $p<0.001$) and memory-conflict resolution ($+0.0941$, $p = 0.005$) improved and no task regressed. Among the comparisons between fine-tuned models, this is the one in which the should-decline difference excludes zero, and it confounds the same three changes as H9.

\paragraph{The traces (descriptive).} All 349 v0.0.2 completions carried a reasoning trace and every trace closed; the traces average 549 characters (median 540, maximum 840) against 575 in the validation targets. Completions averaged 293.5 tokens (median 285, 95th percentile 413) and took 53.8\,s on average on the two T4s (95th percentile 76.3\,s). Six completions reached the 1,024-token budget after a closed trace, all on answerable memory-conflict items, where the answer repeated a phrase; they were scored as generated, at 0.59 to 0.65 each. \code{format\_valid} is 0.0 for all three models.

\par\medskip\noindent\fbox{\parbox{\dimexpr\linewidth-2\fboxsep-2\fboxrule\relax}{\small\textbf{Reported, not decisive.} The secondary rows in this section are uncorrected for multiple comparisons. Three sit near the conventional threshold: memory-conflict resolution under H9 ($p = 0.048$ by group), the H9 should-decline row ($p = 0.054$) and the v0.0.1 to v0.0.2 should-decline row ($p = 0.046$). None of them bears on the pre-registered verdicts, which rest on the primary rows alone.}}

% source: H8b EXP:545-576, L1:429-437; H9 EXP:748-791, L2:165-176; v0.0.1 to v0.0.2 L2:204-210.
\begin{table}[t]
\centering\scriptsize
\caption{Comparisons between fine-tuned models. Every interval is a paired cluster-bootstrap 95\% interval by group; per-task rows give the difference and the verdict by group, with the item-level verdict where it differs. Consistency is the agreement rate by group. Per-task intervals are not recorded; ``n.s.'' is the recorded verdict that the group interval included zero. Sources: \code{docs/experiments.md} H8b and H9 results; Logos v0.0.1 and v0.0.2 run reports.}
\label{tab:head-to-head}
\begin{tabular}{L{27mm}L{40mm}L{40mm}L{40mm}}
\toprule
Row & H8b: Logos v0.0.1 $-$ Hermes & H9: Logos v0.0.2 $-$ Hermes & Logos v0.0.2 $-$ v0.0.1 (secondary) \\
\midrule
Answerable, 271 items, 61 groups (primary) & $-0.0168$ \ci{-0.0546}{+0.0177}; inconclusive & $+0.0409$ \ci{+0.0040}{+0.0780}, $p = 0.031$; better & $+0.0577$ \ci{+0.0234}{+0.0952}, $p<0.001$ \\
Should-decline, 78 items, 17 groups & $-0.0111$ \ci{-0.0429}{+0.0195} & $+0.1020$ \ci{-0.0018}{+0.2062}, $p = 0.054$; not significant & $+0.1131$ \ci{+0.0028}{+0.2296}, $p = 0.046$ \\
Per task, by group & recommendation $-0.1040$ regressed ($p = 0.006$); workspace $-0.0142$ regressed ($p = 0.014$; n.s.\ by item); notification $+0.0837$ n.s.\ (improved by item, $p = 0.007$); briefing $-0.0070$, tool routing $-0.0178$, memory $-0.0204$ n.s.; context not estimable & notification $+0.1324$ ($p<0.001$), recommendation $+0.0893$ ($p = 0.002$), memory $+0.0737$ ($p = 0.048$) improved; tool routing $+0.0330$, workspace $+0.0172$, briefing $-0.0153$ n.s.; context $-0.0002$ not estimable & recommendation $+0.1933$ ($p<0.001$), memory $+0.0941$ ($p = 0.005$) improved; no task regressed \\
Consistency by group, all / answerable & 0.769 / 0.770 $\rightarrow$ 0.833 / 0.853 & 0.769 / 0.770 $\rightarrow$ 0.833 / 0.869 & 0.833 / 0.853 $\rightarrow$ 0.833 / 0.869 \\
Overall mean (descriptive) & 0.8051 vs 0.7896 & 0.8051 vs 0.8596 & 0.7896 vs 0.8596 \\
\bottomrule
\end{tabular}
\end{table}

% source: Hermes HER:310-320; Logos v0.0.1 L1:343-352.
\begin{table}[t]
\centering\small
\caption{Grader sub-scores on the full benchmark. The score is the mean of the ranking, deciding-factor and confidence components where the reference supports them; \code{format\_valid} is reported and excluded. Sources: Hermes and Logos v0.0.1 run reports.}
\label{tab:subscores}
\begin{tabular}{lrrrr}
\toprule
Sub-score & Hermes base & Hermes fine-tuned & Logos v0.0.1 base & Logos v0.0.1 fine-tuned \\
\midrule
score (\code{judgment}) & 0.4755 & 0.8051 & 0.4469 & 0.7896 \\
ranking & 0.5985 & 0.8822 & 0.4468 & 0.8529 \\
\quad nDCG & 0.6174 & 0.9527 & 0.4505 & 0.9459 \\
\quad top-1 accuracy & 0.5702 & 0.7765 & 0.4413 & 0.7135 \\
\quad Kendall's $\tau$ & 0.5449 & 0.8479 & 0.5350 & 0.8345 \\
deciding factor & 0.0860 & 0.6132 & 0.1805 & 0.5989 \\
confidence & 0.7421 & 0.9198 & 0.7135 & 0.9169 \\
\code{format\_valid} & 0.0000 & 0.0000 & 0.0000 & 0.0000 \\
parse failures & 0 & 0 & 0 & 0 \\
\bottomrule
\end{tabular}
\end{table}

\section{Analysis}\label{sec:analysis}

Everything in this section is exploratory: it reads the recorded sub-scores, consistency and abstention tables for what the models learned, and no claim here was registered in advance.

\paragraph{Where the gain comes from.}
% source: HER:310-320, 363-367; L1:343-352; LMD:323-324; SPEC:39; EXP:911.
The sub-scores in Table~\ref{tab:subscores} decompose the H1 gain. For Hermes the deciding-factor component moved most, from 0.0860 to 0.6132: the base model rarely named the factor the policy weighs, the adapter usually does. The ranking component moved from 0.5985 to 0.8822, with nDCG (0.6174 to 0.9527) rising more than top-1 accuracy (0.5702 to 0.7765), one reading is that the adapter orders the tail of a ranking better than it picks its head. Confidence, the component that scores whether the answer commits or declines when the reference does, rose from 0.7421 to 0.9198. Two ceilings bound these figures for every model trained on v0.0.6: because the 78 should-decline items expect labels absent from the training targets (deviation D3), the deciding factor cannot exceed 0.7765 and the score cannot exceed 0.9255 on the full benchmark. On the 271 answerable items the fine-tuned deciding factor is 0.7897 for Hermes (214 correct) and 0.816 for the pilot (221 correct), against 0.1107 (30 correct) for the Hermes base. Logos v0.0.1 shows the same shape with a slightly weaker deciding factor (0.5989) and ranking (0.8529).

\paragraph{Consistency.}
% source: HER:286-292, 559-565; EXP:291-295; HMD:135-136; EXP:573-574, 610-611, 787-788.
Table~\ref{tab:consistency} reports how often a model gives the same decision across the items of one group. Under the configured unit, the scenario family, the fine-tuned Hermes agrees with itself on 0.333 of the 15 test families and is right when it agrees on 0.333, against 0.133 and 0.000 for its base; the pilot reproduces both fine-tuned figures exactly. Those numbers sit under a ceiling of 0.600, because 6 of the 15 families mix different correct labels, which is why the record re-reports consistency by perturbation group, where a perfect model scores 1.000. Hermes rises from 0.346 to 0.769 by group, Logos v0.0.1 from 0.154 to 0.833, and Logos v0.0.2 reaches 0.833 over all groups and 0.869 over the answerable ones. Fine-tuning reduced flips under every perturbation kind for Hermes (paraphrase 53 to 44, evidence order 36 to 21, irrelevant context 24 to 21, length 5 to 3, context order 1 to 0), but paraphrase remains the largest source of inconsistency in both arms. H3's verdict, ``improved, still poor,'' summarizes this: two thirds of the test families (10 of 15) still flip under a perturbation that should not change the decision.

\paragraph{Abstention was learned by scenario family.}
% source: HER:341-354; EXP:306-316; HER:353-354; LMD:500-502.
The lower block of Table~\ref{tab:consistency} splits the should-decline items by the family that produces them. Two families decline unconditionally in training (every one of their 32 and 45 training targets abstains), and on their 50 test items the fine-tuned Hermes declines 50 times out of 50, as does the pilot. Two families decline conditionally, on the evidence (25 of 55 and 3 of 9 training targets abstain), and on their 28 test items both adapters decline 0 times. On the 271 answerable items both adapters commit 271 times. The record's reading is that the models learned ``this kind of question $\rightarrow$ decline,'' not ``the evidence is insufficient $\rightarrow$ decline.'' The base model has the opposite failure: it over-commits almost everywhere, declining correctly on only 3 of the 78 should-decline items. Logos v0.0.2 raised the should-decline score to 0.5855 after training on targets that include the four decline labels, but whether it declines on the evidence rather than on the family was not analyzed.

% source: HER:286-292, 341-347; EXP:291-316; HMD:135-136.
\begin{table}[t]
\centering\scriptsize
\caption{Consistency and abstention. Upper block: agreement across the items of one group, by the configured unit (15 scenario families, where a perfect model scores 0.600) and by perturbation group (78 groups, perfect score 1.000); flips are counts of decision changes attributed to each perturbation kind. Lower block: correct declines per abstention family, and correct commitments on the answerable items. ``Unspecified'' is the record's label for flips not attributed to a named perturbation kind; the record has no row for the schema kind. Sources: Hermes run report; \code{docs/experiments.md} H3.}
\label{tab:consistency}
\begin{tabular}{L{62mm}R{24mm}R{24mm}R{24mm}}
\toprule
Measure & Hermes base & Hermes fine-tuned & Pilot fine-tuned \\
\midrule
Agreement rate by family (ceiling 0.600) & 0.133 & 0.333 & 0.333 \\
Correct-agreement rate by family & 0.000 & 0.333 & 0.333 \\
Families flipping under some perturbation (of 15) & 13 & 10 & 10 \\
Flips: paraphrase / unspecified / evidence order / irrelevant context / length / context order & 53 / 39 / 36 / 24 / 5 / 1 & 44 / 28 / 21 / 21 / 3 / 0 & not reported \\
Agreement rate by group (ceiling 1.000) & 0.346 & 0.769 & not reported \\
\midrule
\code{mem.stale\_explicit\_conflict} (training: 32 of 32 decline), 20 test items & 0 / 20 & 20 / 20 & 20 / 20 \\
\code{route.ask\_when\_underspecified} (45 of 45 decline), 30 test items & 0 / 30 & 30 / 30 & 30 / 30 \\
\code{rec.verify\_before\_recommending} (25 of 55 decline, conditional), 25 test items & 3 / 25 & 0 / 25 & 0 / 25 \\
\code{rec.abstain\_without\_evidence} (3 of 9 decline, conditional), 3 test items & 0 / 3 & 0 / 3 & 0 / 3 \\
Should commit, 271 test items & 256 / 271 & 271 / 271 & 271 / 271 \\
\bottomrule
\end{tabular}
\end{table}

\paragraph{Selection, memorization and the trace.}
% source: HER:146-150; L1:190-196; EXP:710-715; L2:92.
Section~\ref{sec:models} showed that in the two runs with recorded curves the validation loss turns upward in the third epoch while the training loss keeps falling, to about a tenth (Hermes) or a fifteenth (Logos v0.0.1) of the validation loss; the other two runs record only that their final validation loss exceeds the selected checkpoint's. The shipped adapters are the ones selected before that turn, at epochs 1.95 and 1.70, and the benchmark results above belong to those checkpoints, not to the end of training. For Logos v0.0.2 the selection criterion is a loss over the trace and the answer together, and the trace is 58\% of the validation targets by characters; the record keeps no answer-only loss, so it is not known whether the step that minimized the joint loss also minimized the answer loss.

\paragraph{What H9 does and does not say.}
% source: EXP:718-739; LMD:487-490; RM:98; L2:55-57; EXP:733-736.
Logos v0.0.2 differs from Hermes in three ways at once: the base checkpoint (Ministral 3 14B Reasoning instead of Mistral-Nemo 12B), the training data (v0.0.7, with reasoning traces, instead of v0.0.6) and the generation budget (1,024 instead of 512 tokens). The pre-registration counted the first two as the confounded changes and listed the thinking budget as a further one, and the hardware and the torch build differ as well. Hermes was not retrained on v0.0.7, and v0.0.2 has no base arm. H9 therefore says that the new model is measurably better than the old on the answerable items, with a lower bound of $+0.0040$ that the record calls close to zero, and it does not say which change produced the gain or how large each part is. Within the Logos line, v0.0.2 against v0.0.1 isolates the data and budget changes from the base change only partly, since the base changed too (Instruct to Reasoning). The experiment that would separate them, Hermes retrained on v0.0.7 beside a v0.0.2 base arm, is the obvious next run; neither has been made.

\paragraph{Heuristics reported beside their floors.}
% source: HER:299-305; EXP:167-170; L1:333-338; EVAL:261-268.
The faithfulness heuristics moved in the expected direction, and the record declines to make anything of it. For Hermes, citation precision rose from 0.5938 to 0.7192 and the unsupported-claim rate fell from 0.3400 to 0.1600; the number of responses flagged as containing a fabricated citation fell from 147 to 98 (194 to 116 for the pilot, 216 to 94 for Logos v0.0.1). The same heuristic flags 167 of the 349 gold test answers (47.9\%), so a flag is not evidence of a fabrication, and evidence coverage reads 1.0000 for every arm because no benchmark item names decisive evidence identifiers. These numbers are kept in the record as uncalibrated heuristics and are not findings.

\paragraph{Generation economy.}
% source: HER:261-262; L1:293-300; L2:137-144.
Fine-tuning made the answers shorter and faster without any instruction to do so: mean completion length fell from 226 to 94 tokens for Hermes and from 348 to 95 for Logos v0.0.1, and mean latency on the T4 from 24.7 to 16.1\,s and from 36.7 to 17.2\,s, consistent with the adapters having learned the terse format of the targets. The reasoning model reverses part of that: Logos v0.0.2's completions average 293.5 tokens, of which the trace is the larger part, and 53.8\,s on the two-GPU evaluation host; Section~\ref{sec:serving} gives the served latencies.

\section{Serving as a product component}\label{sec:serving}

\paragraph{Packages and pins.}
% source: SVR:12-18, 34; DEP:178-186; KMC/deployment/kleos_logos_v002.yaml:19-23; HER:483-488 (H-F1).
Each served model is a deployment package: the adapter weights, the tokenizer files and a manifest that pins the base checkpoint to a commit and records the adapter's SHA-256 (\code{dc121fa3\ldots} for Hermes, the frozen step-200 weights; \code{e49724f6\ldots} for Logos v0.0.2, the step-175 weights). The research copy of Hermes' adapter configuration recorded no revision (finding H-F1); the package copy pins \code{04d8a905\ldots}, and every later adapter records its revision at training time. Packages live in private model repositories and are read by the serving host at startup. The serving code loads the base in NF4 with float16 compute, applies the adapter and decodes greedily with the evaluated budget (512 tokens for Hermes, 1,024 for Logos v0.0.2), so that the served model is the evaluated model and its outputs can be checked against the frozen evaluation byte for byte. The reply carries the answer (\code{text}) and, for Logos v0.0.2 under the version-2 contract, the trace (\code{reasoning}), which the \kleos{} application shows under ``How Logos decided.''

\paragraph{The free GPU host.}
% source: SVR:89-98; DEP:182, 781-806; HMD:186-188; LMD:467-469; huggingface2026zerogpu.
Both models run on Hugging Face Spaces ZeroGPU \citep{huggingface2026zerogpu}, which allocates half of an NVIDIA RTX Pro 6000 Blackwell (48\,GB, compute capability 12.0) for the duration of a call. A free account has 5 minutes of GPU time per day, charged to the calling account and reset 24 hours after that account's first use; a call is admitted only while 1.5 times its requested duration remains, and it waits up to 60\,s for a GPU. Each call requests
\begin{equation}
\min\bigl(\max(\lceil 10 + \text{max\_new\_tokens} / 12 \rceil,\, 15),\, 60\bigr) \text{ seconds},
\end{equation}
with the constants exposed as Space variables: 53\,s for Hermes and 60\,s for Logos v0.0.2 (the rule gives 96\,s at 1,024 tokens, capped at 60), so admission needs 80\,s and 90\,s of remaining quota respectively. The record estimates roughly 20 to 30 Hermes answers a day and about 9 Logos v0.0.2 answers per fresh daily window; raising its cap to request the full budget's worth of time would cut that to about 5 (all three estimated). A Logos v0.0.2 answer that loops runs out of its 60\,s, is aborted and is reported to the application as a model error, which the application treats like a cut answer.

\paragraph{Reproduction checks.}
% source: Hermes SVR:12-18, 34, 70-71, 102-137; Logos SVR:143-208.
Hermes was checked twice on 2026-09-23. On a Colab T4 the package reproduced 9 of 9 frozen evaluation outputs byte for byte (one prompt from each of the seven tasks plus two should-decline items); the record calls this a reproducibility result, not a score. On ZeroGPU, on different hardware, the same 9 of 9 reproduced byte for byte, 9 of 9 prompt token counts matched and a repeated call returned an identical answer. Logos v0.0.2 was deployed on 2026-10-07 and checked over two daily quota windows, because the first day's window, already partly used before the test, admitted only 7 of the 9 prompts: on day 1 the quota admitted 7 of the 9 prompts, of which 6 reproduced exactly, with prompt tokens equal on all 7; on day 2 the remaining 2 reproduced exactly. In total 8 of 9 answers and traces are identical, prompt tokens match on 9 of 9 and the decisions agree on 8 of 9 (Table~\ref{tab:serving}). The divergent item is an answerable briefing case. Its trace is identical for 568 of its 761 characters; then one sentence goes another way (``Startup has\ldots'' where the evaluation wrote ``Startup is the active workspace\ldots''). The answer declines and asks before looking in another workspace (label \code{ask\_before\_crossing}), where the evaluation had ranked the three items in the active workspace (label \code{scope}, scored 1.0). The record's diagnosis is GPU arithmetic, not setup: the same weights, prompt and decoding rule on a different GPU with float16 rounding. The recorded decision (2026-10-07) was to accept and document: Logos v0.0.2 is served as a Beta, and claims about it say that it measured better than Hermes on the benchmark, not that the live Space reproduces that score. Re-measuring the full benchmark on the Space would take over five weeks at about 9 answers a day (estimated). Appendix~\ref{app:serving} keeps the host facts and the day-by-day smoke record.

% source: SVR:102-116, 131-137, 150-152, 178-195; DEP:793-794; HMD:186-188; LMD:467-469.
\begin{table}[t]
\centering\small
\caption{Serving measurements on ZeroGPU. Hermes figures are medians over 10 calls on 2026-09-23 with answers of 69 to 128 tokens; Logos figures are from the two-day smoke test of 2026-10-07. ``Estimated'' rows follow the admission rule, not a measurement. Source: \code{docs/experiments/serving-verification-records.md}.}
\label{tab:serving}
\begin{tabular}{L{56mm}L{44mm}L{52mm}}
\toprule
Measure & \hermes{} v0.0.6 & \logos{} v0.0.2 \\
\midrule
Generation speed & about 12 tokens/s, cold and warm & 11.0 tokens/s warm, 10.9 cold \\
Warm call, median & 7.9\,s (0.15\,s to obtain a GPU) & 19.7\,s (19.4\,s generating) \\
Cold call, median & 12.9\,s (4.1\,s to obtain a GPU); slowest 34\,s & 29.9\,s (5.5\,s to obtain a GPU) \\
GPU seconds per answer & 7 to 13 & 24 to 26 (217.8 to 232.3 for 9 answers) \\
Peak GPU memory & 8.34\,GiB & 8.86\,GiB \\
Startup: package download / base load and quantization & 2\,s / 41\,s (363 weights) & 1.9\,s / 57.1\,s \\
Reproduction of the frozen evaluation & 9 of 9 (Colab T4) and 9 of 9 (ZeroGPU), byte for byte & 8 of 9 over two quota windows; prompt tokens 9 of 9 \\
Requested GPU duration / quota needed for admission & 53\,s / 80\,s & 60\,s / 90\,s \\
Answers per day on a free account (estimated) & 20 to 30 & about 9 per fresh window \\
\bottomrule
\end{tabular}
\end{table}

\paragraph{Container.}
% source: SVR:39-61.
A Docker image for Hermes was built (6.31\,GB, \code{linux/amd64}) from an allowlisted build context of 57 files, and 257 CPU-side tests passed inside it; it has not been run on a GPU host and is not part of the verified serving path.

\section{Limitations}\label{sec:limitations}

% source: RM:97, 244; EXP:527-529, 731-732; HER:384-392.
\paragraph{One training run per configuration.}
Each adapter was trained once, with seed 42, because of compute limits. The intervals in this paper therefore cover the variability of sampling evaluation items and groups; they do not cover training variance from initialization and data order, which is substantial in fine-tuning \citep{dodge2020finetuning}. Differences between configurations are conditional on one run each and may not replicate across seeds. The agreement between the pilot and Hermes (within 0.004 overall) is consistency across two bases, not a replication of one configuration.

% source: EXP:737-739; RM:157; TDR:153-154; PRV:91-94; V7:67-69.
\paragraph{The data and the benchmark are generated by the same policy.}
Answers, traces and references are all computed from the policies described in Section~\ref{sec:data}, so the benchmark measures how well a model learned those policies, not whether the policies are good or whether the model reasons in general. Test scenarios come from the same families as the training scenarios, with only the input format held out (deviation D2), so the results speak to transfer across format, not to new scenarios. No completed human review exists for v0.0.6 or v0.0.7; the only calibration evidence for the synthetic labels is v0.0.5's review, which rejected 32 of 127 sampled examples before the policy changes that v0.0.6 made, and 48 examples carry a known labelling quirk.

% source: EVAL:215-217; EXP:561-562, 765-766; RM:84; HER:577-583; cameron2008bootstrap.
\paragraph{Few clusters and uncorrected rows.}
The cluster intervals rest on 61 answerable and 17 should-decline groups, and cluster inference with few clusters tends to over-reject. The 8-item context-prioritization task is not estimable by group. Per-task rows are uncorrected for multiple comparisons, and H1's per-task verdicts were never recomputed by group; the protocol amendment left them under their original item-level rule.

% source: EXP:260-265; RM:99, 247.
\paragraph{No model produced the held-out format.}
\code{format\_valid} is 0.0 for every model and arm. Format is excluded from the score by design, so the results measure judgment on prose answers given JSON prompts; whether a model can also produce the JSON structure the product's parsers would prefer is unmeasured.

% source: EXP:313-315; LMD:500-502; EXP:502-504.
\paragraph{Abstention is a family shortcut, and the decline labels were unlearnable for three of the four runs.}
Section~\ref{sec:analysis} shows that the v0.0.6 adapters decline by scenario family, not by evidence, and whether Logos v0.0.2 does better was not analyzed. For the v0.0.6 runs, 78 of 349 items expect labels absent from training (deviation D3), which caps their scores and is why the cross-model comparisons use the answerable subset. Beyond the decline labels, two families are labelled \code{impact} in the test split and only \code{evidence} in training, so 94 of the 349 test items (26.9\%) need a family--label pairing that training never shows.

% source: EXP:445, 449-451, 674-676; RM:251.
\paragraph{Arms not run.}
The bare base model and the fine-tuned-plus-orchestration arm were never evaluated in full, so H7 is only partly measured, and the comparisons use the orchestrated baseline rather than the registered bare one (deviation D1). Logos v0.0.2 has no base arm, so the effect of fine-tuning on the Reasoning checkpoint is not measured; H9 compares two fine-tuned models only.

% source: EXP:718-739; LMD:487-490; RM:98.
\paragraph{H9 confounds three changes.}
The base checkpoint, the training data and the generation budget changed together, the hardware and torch build differ, and Hermes was not retrained on v0.0.7. The interval's lower bound, $+0.0040$, is close to zero and rests on 61 clusters, where cluster inference tends to over-reject; the verdict is one run, one seed, one benchmark.

% source: EXP:727-730, 710-715, 785-786; L1:321-327; HER:391-392; L1:195-196.
\paragraph{Decoding, selection and the base's budget.}
The Reasoning model was decoded greedily although its card recommends sampling at temperature 1; 6 of 349 completions looped and were scored as generated. Its adapter was selected on a loss that covers the trace, and no answer-only loss exists. For Logos v0.0.1, 41 of the base's answers were cut at 512 tokens, which may inflate that run's base-to-adapter gains. Validation losses are not comparable across tokenizers, so the curves in Figure~\ref{fig:validation} compare in shape only.

% source: RM:250; LMD:509-511; SVR:197-208; EXP:939-968; D12 EXP:920; EXP:733-736; TRN:43-44; RM:254; RM:252.
\paragraph{Serving, provenance and tuning.}
The served GPU differs from the evaluated one; Logos reproduced 8 of 9 checks and its live outputs have not been re-measured on the benchmark. The pilot's base revision is not recoverable, so that run is reproducible only modulo upstream not having moved; the Kaggle image was not pinned (deviation D12) and its torch version is recorded only privately. The hyperparameters are engineering defaults, untuned. The faithfulness measures are text heuristics, not entailment checks. A 24B vision-language configuration in the repository was never trained.

\subsection{Data provenance and licensing}
% source: PRV:42-44; RM:334; LMD:35; EXP:91; RM:230; TPE:20-25.
Every training and test example is synthetic, generated from a catalog of fictional entities and captured through an offline mock backend; no release contains conversations, records or other data from real users, and the pipeline's provenance gate rejects any example from a production-observation lane. The two production bases, Mistral-Nemo-Instruct-2407 and Ministral-3-14B-Reasoning-2512, are released under Apache-2.0; the pilot's base, Ministral-8B-Instruct-2410, is under the Mistral AI Research License, which is why that adapter is research-only and was never served. The trained adapters are distributed privately through the deployment packages rather than published; the training and evaluation code, the frozen records and the data pipeline are public under MIT. The misuse surface of a model that ranks a student's own items and declines on weak evidence is narrow, and the serving endpoints require a key and run under a per-account quota.

\section{Reproducibility statement}\label{sec:repro}
% source: RM:14, 174, 228, 230; TDR:267-274; RP:190-196; EXP:939-968; XR:99-132; KCL:46-52; AUD:61-70; HER:57-65; L1:62-68; L2:67-70.
Anyone can run the training and evaluation harness and its 1,638 collected tests without a GPU, evaluate the pipeline end to end on the public development fixtures with an echo backend, and check every configuration hash in this paper except the pilot's against the committed configurations. Anyone can rebuild the data: at the recorded commit of the companion repository, the public pipeline reproduces release v0.0.6 byte for byte (content hash \code{3cc9a744\ldots}) and, from the reasoning change on, release v0.0.7 (\code{b53afa42\ldots}); the test split and the benchmark then carry the SHA-256s cited in Section~\ref{sec:data}. Holders of a release can retrain and re-evaluate each adapter from the runbooks in the repository with the recorded configurations, commits and seeds, and compare their adapters against the recorded SHA-256s and their results against the frozen result-file hashes (Appendix~\ref{app:artifacts}). Nobody can reproduce the pilot's base revision, which was never recorded (deviation D7); its run is reproducible only if the upstream \code{main} pointer has not moved. Every run's manifest records its configuration hash, git commit, dataset hash and runtime adjustments (none), and each verdict in this paper is traceable to a dated entry in the research record. The repository's history was rewritten before publication to remove assistant co-author trailers from seven commit messages; the content, authors and dates of the commits are unchanged, and a table in the record maps each identifier cited by the frozen records to its public counterpart, six of which appear in Appendix~\ref{app:artifacts}.

\section{Conclusion}
A small open-weight model can learn to reproduce a product's decision policy on its benchmark from 820 synthetic examples on a free GPU. The learning is large, visible in every task and consistent across three base checkpoints: fine-tuning moved Hermes from 0.4755 to 0.8051 on a benchmark whose answers are computed by the policy, with a 95\% interval of 0.3068 to 0.3521 on the difference, replicating the pilot, and it moved Logos v0.0.1 by a similar amount. A stronger base alone did not produce a measurably better model; the Reasoning checkpoint of that base, trained on data that adds policy-derived reasoning traces, did, by $+0.0409$ on the answerable items with an interval that excludes zero but starts close to it. What the models learned is specific: they order and label the way the policy does, they decline by the kind of question rather than by the evidence, and they never produce the held-out format. Each model is one run, the second comparison confounds three changes, and the benchmark measures transfer across format within the scenarios the policy was written for. The contribution this paper claims is therefore not a method but an account: a pre-registered, fully logged record of what a student can measure about fine-tuning with the compute a student has, in which the per-task regressions, the inconclusive result and the measurably better one are reported to the same standard.


\section*{Acknowledgements}
Mistral AI released the open-weight base models. Hugging Face maintains Transformers, PEFT and bitsandbytes and hosts the models on Spaces ZeroGPU. Google Colab and Kaggle provided the free GPU sessions on which every run was trained and evaluated.

\bibliographystyle{plainnat}
\bibliography{references}

\clearpage
\appendix
% Appendices. Every value is copied from the record named in the "% source:" comment above it.

\section{Hyperparameters and software}\label{app:hyper}
Table~\ref{tab:hyper-full} lists the settings that Table~\ref{tab:recipe} omits and the software versions of each run.
% source: TRN:46-62; AUD:47-53, 72, 29; HER:159-161, 213, 498; L1:252; KMC/training/*.yaml; versions HER:67-68, L1:70-72, L2:72, EXP:875, AUD:74, SVR:74, 151.
\begin{table}[H]
\centering\small
\caption{Training settings not shown in Table~\ref{tab:recipe}, and the software versions per run. ``not recorded'' means the public record keeps no value. Sources: \code{docs/training.md}, the training configurations under \code{configs/training/} and the run reports.}
\label{tab:hyper-full}
\begin{tabular}{L{52mm}L{100mm}}
\toprule
Setting & Value \\
\midrule
Logging interval & every 5 steps (all four configurations) \\
Validation and checkpoint interval & 50 steps (pilot, Hermes); 25 steps (Logos v0.0.1, v0.0.2) \\
Checkpoints retained & 3, plus the best (\code{load\_best\_model\_at\_end} on the lowest validation loss) \\
Strict configuration checking & off (pilot, Hermes); on (Logos v0.0.1, v0.0.2) \\
PEFT options & \code{use\_dora} false, \code{use\_rslora} false, \code{modules\_to\_save} none, default LoRA initialization (recorded for the pilot and Hermes; not inspected for Logos v0.0.1) \\
Warmup & ratio 0.03, resolved by Transformers 5 to 10 steps \\
Mistral tokenizer regex fix & off for Hermes; on for both Logos runs (changes 0 of 1,350 examples) \\
Adapter files (fp32) & pilot 174,655,536 bytes; Hermes 228,140,600; Logos v0.0.1 243,869,280; Logos v0.0.2 not recorded \\
Trainable share reported by PEFT & pilot 1.5293\%; Hermes 1.38\% \\
fp16-skipped optimizer steps & Hermes 4 (through step 305); Logos v0.0.1 2 (both before step 250); pilot and v0.0.2 not recorded \\
\midrule
Pilot software & transformers 5.16.1, peft 0.20.0, bitsandbytes 0.50.2, torch 2.11.0+cu128, Python 3.13.15 \\
Hermes software & torch 2.11.0+cu128, transformers 5.16.1, peft 0.20.0, accelerate 1.14.0, bitsandbytes 0.50.2, datasets 4.8.5 \\
Logos v0.0.1 software & as Hermes, plus tokenizers 0.23.2 \\
Logos v0.0.2 software & transformers 5.16.1, peft 0.20.0, accelerate 1.14.0, bitsandbytes 0.50.2, tokenizers 0.23.2 installed over Kaggle's image; torch is Kaggle's own, recorded only in the run's environment file \\
Serving (Logos Space) & torch 2.11.0+cu128, transformers 5.16.1, peft 0.20.0, bitsandbytes 0.50.2, gradio 6.28.0, spaces 0.51.3; the Hermes Space pins pydantic 2.12.5 \\
\bottomrule
\end{tabular}
\end{table}

\section{Prompt and answer format}\label{app:format}
% source: data/examples/synthetic_train.jsonl (fixture-prior-0001, public development fixture); KMC/models/ministral3_14b_reasoning.yaml:33-37; V7:93-119.
An example is a conversation of three messages. The sealed releases are not published as files (the public pipeline rebuilds them byte for byte), so the format is shown here on the first public development fixture shipped with the training harness, which its metadata marks as a synthetic development fixture and not the research dataset. It has the same structure as a release example without a reasoning trace.

\begin{small}
\begin{verbatim}
system:  You are the reasoning layer of KLEOS. Deterministic systems have already
         retrieved and ranked the context. Decide, then justify from the evidence
         provided. Do not invent facts.

user:    Rank these items by what deserves attention first.

         - internship-application: deadline closes in 2 days; evidence strong
           (confirmed opening)
         - networking-event: deadline in 3 weeks; evidence weak (no confirmed
           attendees)
         - resume-review: deadline no deadline; evidence medium (peer feedback
           offered)

assistant:
         1. internship-application
         2. resume-review
         3. networking-event

         Reasoning: internship-application comes first because its deadline is
         the nearest (closes in 2 days) and the evidence of impact is strong
         (confirmed opening). Items with weak evidence or no deadline rank last,
         since acting on them displaces work with a firm cost of delay.
\end{verbatim}
\end{small}

In release v0.0.7 every training and validation answer ends with a line ``What decided it: $\langle$label$\rangle$.'' and carries a \code{reasoning} field. For the Reasoning base the chat template renders that field before the answer, so a supervised completion is \code{[THINK]}trace\code{[/THINK]}answer\code{</s>}, with \code{[THINK]} as token 34 and \code{[/THINK]} as token 35, and the loss covers the trace and the answer. The trace follows five steps:
\begin{enumerate}
  \item the rule the family's policy weighs, for example ``Score each item as deadline $\times$ evidence $\times$ impact'';
  \item one line per item, in prompt order, mapping the prompt's wording to the policy's values, for example ``Redpine: due in 2 days (0.33), reported (0.55), high (1.0) $\rightarrow$ 0.18'';
  \item each check the policy applied and its outcome: close call (relative separation below 0.15), stale explicit statement (30 days or older), usable evidence (weight 0.55 or above), workspace scope, strong challenger (weight 0.8 or above) and ambiguous request;
  \item why the deciding factor decided it;
  \item the conclusion: the ranking for a committed answer, or the failed check and the question to ask for a declined one.
\end{enumerate}
Traces are capped at 1,200 characters and follow wording rules so that they pass the privacy and fact checks: only item names are capitalized and never two capitalized words in a row, numbers use at most two decimal places, and no dates, at-signs or comma-grouped numbers appear. Each step has two or three phrasings, chosen deterministically, so a rebuild reproduces every trace.

\section{Policies, labels and the test population}\label{app:policies}
Table~\ref{tab:policies} reproduces the twelve policies and Table~\ref{tab:labels} the eleven deciding-factor labels with the policy that produces each; the scenario families are in Table~\ref{tab:families}.
% source: the policy table of Kleos-Training-Data docs/scenarios.md (SC:194-207), copied cell by cell.
\begin{table}[H]
\centering\scriptsize
\caption{The twelve decision policies: what each orders by, when it declines, and the deciding factor it names. Copied from the policy table of Kleos-Training-Data, \code{docs/scenarios.md}.}
\label{tab:policies}
\begin{tabular}{L{43mm}L{31mm}L{40mm}L{40mm}}
\toprule
Policy & Orders by & Abstains when & Deciding factor \\
\midrule
\code{rank\_by\_deadline\_then\_evidence} & Score & Never & Deadline, evidence or impact, whichever differs most between the top two \\
\code{rank\_or\_abstain\_when\_close} & Score & The top two are within a 15\% relative separation & \code{insufficient\_separation} \\
\code{rank\_by\_reliability\_over\_recency} & Evidence, then recency & Never & Evidence, or deadline when evidence ties \\
\code{resolve\_or\_abstain\_on\_support} & Evidence, then recency & The top two are within 15\% on evidence and were recorded at the same time & \code{insufficient\_separation} \\
\code{select\_by\_evidence\_need} & Evidence $\times$ impact, then freshness & Never & Evidence or impact \\
\code{rank\_by\_relevance\_over\_recency} & Impact (relevance), then recency & Never & Impact, or deadline when impact ties \\
\code{respect\_workspace\_scope} & In-scope items first, then score & Never. Out-of-scope items are flagged, not acted on. & \code{scope} when scopes are mixed, otherwise the score ranking \\
\code{ask\_before\_crossing\_workspace} & Same as \code{respect\_}\allowbreak\code{workspace\_scope} & No in-scope item is adequate and the strongest out-of-scope item is & \code{ask\_before\_crossing} \\
\code{defer\_to\_explicit\_statement} & Explicit statements, then evidence, then recency & The explicit statement is at least 30 days old and a newer inferred record has evidence of at least 0.8 & \code{explicit\_statement}, \code{evidence} or \code{stale\_explicit\_conflict} \\
\code{verify\_when\_evidence\_weak} & Score & No item is adequate & \code{missing\_input} \\
\code{ask\_when\_request\_ambiguous} & Score & The family sets \code{prompt.request\_ambiguous} & \code{request\_ambiguous} \\
\code{prefer\_least\_privilege\_source} & Adequacy, then narrowest reach, then evidence & Never & \code{least\_privilege} when several sources are adequate, otherwise \code{evidence} \\
\bottomrule
\end{tabular}
\end{table}

% source: labels and the policy cells that produce them, SC:194-207 (as consolidated in the fact pack, A6); test counts EXP:909 (D3), V7:130-132.
\begin{table}[H]
\centering\footnotesize
\caption{The eleven deciding-factor labels named by the policy table, the policy-table cell that produces each (quoted verbatim, with the policy named), and what the record says about the test split. The benchmark's label set has nine options; the two labels whose only families have no JSON-format examples are absent from it. Test counts are recorded for the four decline labels only. Source: Kleos-Training-Data, \code{docs/scenarios.md}, policy table; \code{docs/experiments.md}, deviation D3; the v0.0.7 repair specification.}
\label{tab:labels}
\begin{tabular}{L{36mm}L{12mm}L{70mm}L{26mm}}
\toprule
Label & Kind & Policy-table cell & Test split \\
\midrule
\code{deadline} & commit & \code{rank\_by\_deadline\_then\_evidence}: ``Deadline, evidence or impact, whichever differs most between the top two''; \code{rank\_by\_reliability\_over\_recency}: ``Evidence, or deadline when evidence ties''; \code{rank\_by\_relevance\_over\_recency}: ``Impact, or deadline when impact ties'' & count not recorded \\
\code{evidence} & commit & the three cells above; \code{select\_by\_evidence\_need}: ``Evidence or impact''; \code{prefer\_least\_privilege\_source}: ``\code{least\_privilege} when several sources are adequate, otherwise \code{evidence}''; \code{defer\_to\_explicit\_statement}: ``\code{explicit\_statement}, \code{evidence} or \code{stale\_explicit\_conflict}'' & count not recorded \\
\code{impact} & commit & \code{rank\_by\_deadline\_then\_evidence} and \code{select\_by\_evidence\_need} (cells above); \code{rank\_by\_relevance\_over\_recency}: ``Impact, or deadline when impact ties'' & count not recorded \\
\code{scope} & commit & \code{respect\_workspace\_scope}: ``\code{scope} when scopes are mixed, otherwise the score ranking'' & count not recorded \\
\code{explicit\_statement} & commit & \code{defer\_to\_explicit\_statement}: ``\code{explicit\_statement}, \code{evidence} or \code{stale\_explicit\_conflict}'' & count not recorded \\
\code{least\_privilege} & commit & \code{prefer\_least\_privilege\_source}: ``\code{least\_privilege} when several sources are adequate, otherwise \code{evidence}'' & absent: its family has no JSON-format examples \\
\code{insufficient\_separation} & decline & \code{rank\_or\_abstain\_when\_close} abstains when ``The top two are within a 15\% relative separation''; \code{resolve\_or\_abstain\_on\_support} when ``The top two are within 15\% on evidence and were recorded at the same time'' & 3 items \\
\code{ask\_before\_crossing} & decline & \code{ask\_before\_crossing\_workspace} abstains when ``No in-scope item is adequate and the strongest out-of-scope item is'' & absent: its family has no JSON-format examples \\
\code{stale\_explicit\_conflict} & decline & \code{defer\_to\_explicit\_statement} abstains when ``The explicit statement is at least 30 days old and a newer inferred record has evidence of at least 0.8'' & 20 items \\
\code{missing\_input} & decline & \code{verify\_when\_evidence\_weak} abstains when ``No item is adequate'' & 25 items \\
\code{request\_ambiguous} & decline & \code{ask\_when\_request\_ambiguous} abstains when ``The family sets \code{prompt.request\_ambiguous}'' & 30 items \\
\bottomrule
\end{tabular}
\end{table}

\section{Deviations D1 to D13}\label{app:deviations}
Table~\ref{tab:deviations} condenses the deviations log; the full text of each entry is in the research record.
% source: EXP:909-921 (the deviations log), condensed; the recorded effect on the result is quoted or paraphrased from each row.
\begin{table}[H]
\centering\scriptsize
\caption{The deviations log, condensed. Each row's last column is the effect recorded at the time. Source: \code{docs/experiments.md}, deviations log.}
\label{tab:deviations}
\begin{tabular}{L{8mm}L{16mm}L{70mm}L{58mm}}
\toprule
Id & Date & Deviation & Recorded effect \\
\midrule
D1 & 2026-09-15 & H1 registered \code{arm0\_base} vs \code{arm2}; the run compared \code{arm1\_base\_orchestrated} vs \code{arm2} & The research question names the orchestrated baseline; arm0 exists only as a 20-example probe (0.5659), a biased head-of-file slice not comparable with a full run \\
D2 & 2026-09-15 & H1 registered an entity holdout; the sealed release uses a format holdout on JSON & The test split probes format transfer, not entity generalization \\
D3 & 2026-09-15 & 78 of 349 test items (22\%) expect a label that appears zero times in training (\code{request\_ambiguous} 30, \code{missing\_input} 25, \code{stale\_explicit\_conflict} 20, \code{insufficient\_separation} 3) & Affects both arms identically, so H1 stays fair; absolute deciding-factor figures carry the caveat (0.6332 overall, 0.816 on answerable items for the pilot) \\
D4 & 2026-09-15 & The ranking extractor was corrected before the run (it had returned whole prose lines) & Instrument changed before any result; no reported number affected \\
D5 & 2026-09-15 & The formatter folds the system prompt into the first user turn when the template drops it & Fixed before the run; without it every example would have trained without its policy instructions \\
D6 & 2026-09-15 & The composite \code{kleos\_policy} grader was added, with \code{format\_valid} reported separately and excluded & Without the separation a format failure is indistinguishable from a judgment failure \\
D7 & 2026-09-15 & The pilot used \code{revision: main}; the resolved commit is not recoverable & Future runs are pinned; the historical record is not backfilled \\
D8 & 2026-09-18 & Hermes training was interrupted after step 250 and resumed 18 hours later on another T4 & No effect: the selected checkpoint (step 200) predates the interruption; resume state verified (step-300 evaluations agree to $1.4 \times 10^{-5}$) \\
D9 & 2026-09-27 & Logos v0.0.1 trained over four Colab sessions with two resumes (from steps 100 and 250) & No effect on the selected adapter (step 175, written in session 2); the reported \code{train\_loss} is a resume artifact and is not quoted \\
D10 & 2026-09-29 & Logos \code{arm2} was evaluated in three sessions; the first two never reached storage (finding L-F3); the third generated all 349 in one pass & No effect: a single uninterrupted run at the registered commit \\
D11 & 2026-09-29 (declared before it ran) & Logos \code{arm1} would run at a later commit containing the L-F3 file-handling fix & Done as declared on 2026-10-01; the change touches file handling only; the second session resumed 172 recorded generations \\
D12 & 2026-10-06 & Both Kaggle notebooks ran on the ``Latest Container Image,'' not a pinned one & No effect: neither run was resumed in a later version; libraries were installed at pinned versions \\
D13 & 2026-10-06 & The training run is Version 2 of its notebook; Version 1 stopped at the commit guard in its first cell & No effect: Version 1 ran no code beyond the guard; both runs used the registration commit and the configuration hash matched before training \\
\bottomrule
\end{tabular}
\end{table}

\section{Findings}\label{app:findings}
Table~\ref{tab:findings} lists the three findings registers.
% source: AUD:133-224 (F1-F3); HER:483-583 with statuses HMD:222-237 (H-F1 to H-F14); LF:20-118 with statuses LMD:536-545 (L-F1 to L-F8).
\begin{table}[H]
\centering\scriptsize
\caption{The findings registers, one line each, with status as of 2026-10-07. Sources: the artifact audit, the Hermes run report and \code{docs/experiments/logos-findings.md}.}
\label{tab:findings}
\begin{tabular}{L{10mm}L{92mm}L{50mm}}
\toprule
Id & Finding & Status \\
\midrule
F1 & nDCG could exceed 1.0 on degenerate rankings; the pilot baseline re-graded from 0.5231 to 0.4744 (210 of 349 responses changed), the fine-tuned arm unchanged & resolved 2026-09-15 \\
F2 & A 17.1\,MB \code{tokenizer.json} was refused by the 5\,MB scan cap & resolved 2026-09-15 \\
F3 & The base-model revision was not pinned & resolved for future runs 2026-09-15; the pilot is not backfilled (D7) \\
\midrule
H-F1 & The research adapter configuration records \code{revision: null} & research copy kept; the package pins the revision; later adapters record it \\
H-F2 & The manifest contradicts itself on the trainable count (1,342,592,000 against 57,016,320) & open; the per-layer reconciliation gives 57,016,320 \\
H-F3 & A stale sliding-window note was frozen into the records & open \\
H-F4 & A checkpoint-saved event is not a durability guarantee on Drive & open; weightless checkpoints are rejected at resume \\
H-F5 & The comparison tool checked benchmark identity by path, not content & fixed for later runs \\
H-F6 & The comparability check never compares base models & open (methodological) \\
H-F7 & Result files drop the grader details & fixed for later runs \\
H-F8 & \code{train\_loss} is invalid after a resume: the reported value is a resume artifact; the mean over the second session's steps 251 to 309 is 0.00459 & open; never quoted \\
H-F9 & The retention pass did not protect the best checkpoint & fixed \\
H-F10 & The memory estimator was about 5\,GB optimistic (7.96 to 8.12 against 13.09\,GB) & fixed \\
H-F11 & Consistency was grouped by scenario family (ceiling 0.600, 57 flips for a perfect model) & reported by group from H8 on \\
H-F12 & The fabricated-citation heuristic flags gold answers (47.9\% of test targets) & reported beside its floor \\
H-F13 & Evidence coverage is vacuous (every response scores 1.0) & reported as vacuous \\
H-F14 & Item-level intervals treat 349 dependent items as independent & cluster intervals from H8 on \\
\midrule
L-F1 & LoRA scoping for the vision-language adapter fails on Transformers 5 (7 of 14 modules landed in the vision tower) & open; the strict test is an expected failure \\
L-F2 & The tokenizer class depends on \code{config.json} (442 against 492 tokens for the same text) & prompt tokens matched on 9 of 9 served prompts \\
L-F3 & Evaluation resume lost killed sessions' work on the Drive mount & fixed; confirmed 2026-10-01 (172 generations resumed) \\
L-F4 & A fresh start into a folder with checkpoints would delete the best one & open \\
L-F5 & Checkpoint validation ignores optimizer and scheduler state & open \\
L-F6 & \code{device\_map: auto} split the model unevenly across two T4s (probe 10.16\,GB on GPU 1 against an estimated 8.53; run peak 10.21\,GB) & open; no effect on the run \\
L-F7 & The stop reason misses a budget cut after a closed trace (6 of 349) & resolved in serving \\
L-F8 & Evaluation reports no progress while running & open \\
\bottomrule
\end{tabular}
\end{table}

\section{Artifact register}\label{app:artifacts}
% source: config hashes AUD:61, HER:57, L1:62, L2:67; commits AUD:68-69, HER:64, L1:67-68, L2:68 and the map XR:106-132; dataset identity AUD:64, HER:60-61, L1:64, L2:70, KMC/deployment/kleos_logos_v002.yaml:23; adapters AUD:29, HER:213, L1:252, KMC/deployment/kleos_logos_v002.yaml:20; result files HER:462-463, L1:471-472, 492, L2:220; test and benchmark SPEC:54, EXP:477; KCL:46-52.
Three kinds of dataset identity appear in the records and are never interchangeable: the \emph{release content hash}, computed by the data pipeline over a sealed release (\code{3cc9a74486c42e8e} for v0.0.6, \code{b53afa4216bf6973} for v0.0.7, as the changelog prints them); the harness \emph{dataset hash}, computed over the loaded split files (\code{c7340f1d9cb31b1d\allowbreak{}1d7f0ade65d83740\allowbreak{}159a887f5d6cb495\allowbreak{}5820065076ffa875} for v0.0.6); and the deployment record's \emph{dataset SHA-256} (\code{5d02a5afd7be18e5\allowbreak{}7f13b7d86fa8c31a\allowbreak{}1b008dd9e1927f4f\allowbreak{}a506356649699c48} for v0.0.7). Likewise the test split (\code{a4decaaf029b2273\allowbreak{}66846d44018b0b1b\allowbreak{}669d21f068800a9f\allowbreak{}3175e15ec4783980}) and the benchmark built from it (\code{a11ffad75f5147f9\allowbreak{}d0ddad7bad4bfc07\allowbreak{}3dc730df2b19173f\allowbreak{}f642774233b4b266}) are different objects.

Table~\ref{tab:artifacts} gives each run's identity.

\begin{table}[H]
\centering\scriptsize
\caption{Per-run identity. Commits are given as the frozen records cite them and as they appear in the public history after the rewrite described in Section~\ref{sec:repro}; manifests hold the former. The pilot's configuration hash no longer reproduces at the current head because its base was pinned afterwards. Sources: the artifact audit, the run reports, the deployment record and \code{docs/experiments/README.md}.}
\label{tab:artifacts}
\begin{tabular}{L{22mm}L{130mm}}
\toprule
Run & Identity \\
\midrule
Ministral-8B pilot & config hash \code{3fbb3f90ed9662ee\allowbreak{}f801363820ddaade\allowbreak{}0c24e591468447b0\allowbreak{}30ef50d12b7806d3}; commit \code{2ad7c9a} (public \code{b64099b}), clean; dataset hash \code{c7340f1d\ldots}; adapter \code{6d4a0a93d274160b\allowbreak{}aeea65843f48d57d\allowbreak{}f759cddc9d423f55\allowbreak{}a20d7c88659e473c}, 174,655,536 bytes \\
\hermes{} v0.0.6 & config hash \code{b2328857c6026dd7\allowbreak{}c904a6acadd00ec5\allowbreak{}5a6ce5289bf6210e\allowbreak{}743615fce6d8d3e8}; commit \code{4cd76c4} (public \code{22ca978}), clean; dataset hash \code{c7340f1d\ldots}; adapter \code{dc121fa36ce1409c\allowbreak{}a142e8da61a809f9\allowbreak{}386e32d0ef09c421\allowbreak{}4efa271d5b857b32}, 228,140,600 bytes; result files arm1 \code{71dee4c0\ldots}, arm2 \code{428400f5\ldots}, arm2 annotated \code{99fdcdc3\ldots} \\
\logos{} v0.0.1 & config hash \code{18008c6716a58afc\allowbreak{}284c64eb7e1e96c9\allowbreak{}bfce34b8da2b8c48\allowbreak{}9b6d59c0c45b6f71}; training commit \code{ec4f9e3} (public \code{81e3e9a}); arm1 evaluation commit \code{ed5a987} (public \code{e915060}); dataset hash \code{c7340f1d\ldots}; adapter \code{f3e8dcdc09a3f3bb\allowbreak{}9f6fc4f64f867ff8\allowbreak{}b97834c955a03fcb\allowbreak{}de98a8068d5570a7}, 243,869,280 bytes; result files arm1 \code{a0b048b1\ldots}, arm2 \code{09243630\ldots} \\
\logos{} v0.0.2 & config hash \code{d1961583546b5761\allowbreak{}dd854cfe845838ca\allowbreak{}91a87ea6189d54be\allowbreak{}617de21085e8cfa9}; commit \code{a17ace7} (public \code{6d7e450}); release content hash \code{b53afa42\ldots}; adapter \code{e49724f6554db662\allowbreak{}769fbe9b9431105d\allowbreak{}bc12f8f75cecd6e0\allowbreak{}e00728ea3f04dc4f}; result file arm2 \code{0dd74661\ldots}, whose \code{benchmark\_fingerprint} field is \code{c3d0d1f1\ldots} (a hash stored by the harness, distinct from the benchmark SHA-256) \\
\bottomrule
\end{tabular}
\end{table}

The Logos configuration hashes include the literal dataset and output paths used on the training hosts, which the record keeps for that reason. The pre-registration commits map as follows: H5 and H6 on 2026-09-15, \code{7145edd} to \code{a6bae0c}; H8 and the clustered-interval amendment on 2026-09-24, \code{ec4f9e3} to \code{81e3e9a}; D11 on 2026-09-29, \code{ed5a987} to \code{e915060}; H9 on 2026-10-06, \code{a17ace7} to \code{6d7e450}; the first commit, \code{fd87931} of 2026-08-14, is unchanged.

\section{Serving details and the smoke record}\label{app:serving}
Table~\ref{tab:zerogpu} records the host's rules and the checks behind Section~\ref{sec:serving}.
% source: SVR:89-98 (ZeroGPU facts); DEP:182, 781-794 (duration rule); SVR:12-18, 102-137 (Hermes); SVR:150-195 (Logos day 1 and day 2); L2:136-145, 185-194 (thinking statistics).
\begin{table}[H]
\centering\small
\caption{ZeroGPU facts as recorded on 2026-09-23 and the two-day Logos smoke record of 2026-10-07. Source: \code{docs/experiments/serving-verification-records.md}.}
\label{tab:zerogpu}
\begin{tabular}{L{50mm}L{102mm}}
\toprule
Item & Recorded value \\
\midrule
GPU & half an NVIDIA RTX Pro 6000 Blackwell: 48\,GB, compute capability 12.0 \\
Daily quota & 5 minutes of GPU time for a free account, 2 minutes for an unauthenticated caller; charged to the calling account; reset 24 hours after that account's first use \\
Admission and queue & admitted while remaining quota covers 1.5 times the requested duration; waits up to 60\,s for a GPU; up to 2 ZeroGPU Spaces per account; no persistent storage \\
Duration rule & $\min(\max(\lceil \text{BASE} + \text{max\_new\_tokens}/\text{TOKENS\_PER\_SECOND} \rceil, \text{MIN}), \text{MAX})$ with defaults 10\,s, 12 tokens/s, 15\,s and 60\,s; measured speeds of 12 (Hermes) and 11.0 (Logos) tokens/s left the defaults in place \\
Hermes, Colab check & adapter \code{dc121fa3\ldots}; the package copy of the adapter configuration pins \code{04d8a905\ldots}; tokenizer loaded without the regex fix, as trained; 9 of 9 byte-identical \\
Hermes, ZeroGPU (10 calls) & 9 of 9 byte-identical, token counts equal, repeated call identical; 7 of 10 calls got a new worker; 73 to 131\,s of the 300\,s daily quota consumed; startup: download 2\,s, load and quantization 41\,s, packing 8.54\,GB 3\,s; peak CPU memory at startup not recorded \\
Logos, startup & download 1.9\,s, load 57.1\,s, host peak 30.6\,GiB of 2,000\,GiB, 9.20\,GB packed for the GPU \\
Logos, day 1 & 7 of 9 prompts admitted before the quota refused the 8th; 6 of 7 identical; 161 to 174 GPU seconds; warm median 19.7\,s, cold 29.9\,s \\
Logos, day 2 & 2 of 2 identical (a tool-routing and a workspace-reasoning item); cold 25.2\,s (236 tokens), warm 33.6\,s (a 406-token answer); 56.5 to 58.8 GPU seconds \\
Logos, totals & 8 of 9 identical in answer and trace; prompt tokens equal on 9 of 9; decisions agree on 8 of 9; 217.8 to 232.3 GPU seconds (24 to 26 per answer) \\
Logos, evaluation traces & 349 of 349 answers traced; 0 traces left open; 0 empty answers or parse failures; mean 549 characters (median 540, maximum 840) against training traces of 327 to 878; completions mean 293.5 tokens, median 285, 95th percentile 413 \\
\bottomrule
\end{tabular}
\end{table}


\end{document}
